% REMEMBER: You must not plagiarise anything in your report. Be extremely careful.
\documentclass{l4proj}

    
%==============================================================================
% Put any additional packages here
% You can add any packages you want, as long as it does not alter
% the overall format (e.g. don't change the margins or the reference style).
%
\usepackage{pdfpages} % if you want to include a PDF for an ethics checklist, for example
\usepackage{float}
\usepackage{booktabs}
\usepackage{listings}
\usepackage{xcolor}
\usepackage{siunitx}
\usepackage{subcaption}
\usepackage{algorithm2e}
\definecolor{codebg}{rgb}{0.97,0.97,0.97}
\definecolor{codestring}{rgb}{0.13,0.54,0.13}
\definecolor{codekw}{rgb}{0.0,0.0,0.7}
\definecolor{codecomment}{rgb}{0.4,0.4,0.4}

\lstdefinestyle{pythonstyle}{
  backgroundcolor=\color{codebg},
  basicstyle=\ttfamily\small,
  keywordstyle=\color{codekw}\bfseries,
  stringstyle=\color{codestring},
  commentstyle=\color{codecomment}\itshape,
  breaklines=true,
  showstringspaces=false,
  frame=single,
  framesep=3pt,
  xleftmargin=6pt,
  xrightmargin=3pt,
}

\lstdefinestyle{yamlstyle}{
  backgroundcolor=\color{codebg},
  basicstyle=\ttfamily\small,
  keywordstyle=\color{codekw},
  stringstyle=\color{codestring},
  commentstyle=\color{codecomment}\itshape,
  breaklines=true,
  showstringspaces=false,
  frame=single,
  framesep=3pt,
  xleftmargin=6pt,
}

\lstdefinestyle{hoconStyle}{
  backgroundcolor=\color{codebg},
  basicstyle=\ttfamily\small,
  commentstyle=\color{codecomment}\itshape,
  stringstyle=\color{codestring},
  breaklines=true,
  showstringspaces=false,
  frame=single,
  framesep=3pt,
  xleftmargin=6pt,
}

\begin{document}

%==============================================================================
%% METADATA
\title{A Framework for Benchmarking SQL Workloads on Distributed Data Lakehouses} % change this to your title
\author{2736414Y}
\date{September 14, 2026}

\maketitle

%==============================================================================
%% ABSTRACT
\begin{abstract}
    The Apache Trino and Iceberg ecosystem lacks a reproducible, systematic benchmarking framework: existing tools require significant manual configuration, lack integrated hardware monitoring, or cannot port across deployment environments without modification. TriBench is a cross-platform benchmarking framework that addresses this gap by providing declarative experiment definitions, concurrent hardware and per-pod resource monitoring, and persistent result storage in a portable, self-contained bundle. The same experiment specification runs unchanged on a local Docker deployment or a multi-node Kubernetes cluster by switching a single configuration profile. Evaluated using TPC-H and TPC-DS across three environments --- local Docker, bare-metal Kubernetes at the University of Glasgow, and Google Kubernetes Engine --- TriBench produced reproducible query timings with full telemetry at scale factors up to SF~100, fulfilling the identified need for a PEEL-equivalent systematic benchmarking tool in the Trino ecosystem.

    
\end{abstract}

%==============================================================================
%% ACKNOWLEDGEMENTS
\chapter*{Acknowledgements}
% Enter any acknowledgements here. This is optional; you may leave this blank if you wish,
% or remove the entire chapter
%
% We give thanks to the Gods of LaTeX, who in their eternal graciousness, 
% have granted that this document may compile without errors or overfull hboxes.
%

%==============================================================================

% EDUCATION REUSE CONSENT FORM
% If you consent to your project being shown to future students for educational purposes
% then insert your name and the date below to  sign the education use form that appears in the front of the document. 
% You must explicitly give consent if you wish to do so.
% If you sign, your project may be included in the Hall of Fame if it scores particularly highly.
%
%
\def\consentname {Lin Yuan} % your full name
\def\consentdate {20 March 2026} % the date you agree
%
\educationalconsent


%==============================================================================
\tableofcontents

%==============================================================================
%% Notes on formatting
%==============================================================================
% The first page, abstract and table of contents are numbered using Roman numerals and are not
% included in the page count. 
%
% From now on pages are numbered
% using Arabic numerals. Therefore, immediately after the first call to \chapter we need the call
% \pagenumbering{arabic} and this should be called once only in the document. 
%
%
% The first Chapter should then be on page 1. 

% PAGE LIMITS
% You are allowed 40 pages for a 40 credit project and 30 pages for a 
% 20 credit report. 
% This includes everything numbered in Arabic numerals (excluding front matter) up
% to but *excluding the appendices and bibliography*.
%
% FORMATTING
% You must not alter text size (it is currently 10pt) or alter margins or spacing.
% Do not alter the bibliography style. 
%
%==================================================================================================================================
%
% IMPORTANT
% The chapter headings and structure here are **suggestions**. You don't have to follow this model if
% it doesn't fit your project. Every project should have an introduction and conclusion,
% however.  If in doubt, your supervisor can give you specific guidance; their view takes precedence over
% the structure suggested here.
%
%==================================================================================================================================
\chapter{Introduction}
\label{chap:intro}

% reset page numbering. Don't remove this!
\pagenumbering{arabic}

This chapter outlines the motivation for developing TriBench, a cross-platform benchmarking framework for SQL workloads on distributed data lakehouses. The key challenges in the problem space are identified, and the overall aims of the project are introduced.

\section{Motivation}
\label{sec:motivation}

The past decade has seen a significant shift in how organisations manage and query large-scale analytical data. Traditional data warehouses offered strong consistency and mature query optimisation, but their proprietary formats and high storage costs made them impractical for the volumes of raw data modern organisations collect. Data lakes addressed the storage problem by writing data as files on cheap object storage, but at the cost of ACID guarantees, schema enforcement, and query performance. The data lakehouse architecture \citep{armbrust2021lakehouse} emerged to bridge this gap, combining the low-cost object storage of a data lake with the reliability guarantees of a warehouse through an open table-format layer. Apache Iceberg \citep{apache_iceberg} has become the dominant realisation of this approach, providing snapshot isolation, time-travel queries, and schema evolution on top of columnar Parquet files. Apache Trino \citep{sethi2019presto}, a distributed SQL engine originally developed at Meta, is widely deployed as the interactive query layer over Iceberg-based lakehouses in production environments.

Rather than a single system under test, a Trino lakehouse deployment requires the coordinated operation of five interdependent services: Trino itself (coordinator and one or more workers), a Hive Metastore for the Iceberg catalogue, a PostgreSQL backend for catalogue persistence, and an S3-compatible object store such as MinIO for table data. Each service carries its own hardware-dependent configuration parameters --- JVM heap sizes, memory limits, worker counts, and network addresses --- that must be set consistently before any experiment can run. Changing a single parameter, such as increasing the number of Trino workers to study scale-out behaviour, requires coordinated changes across multiple configuration files and a fresh deployment of the affected services. Running such an experiment by hand is error-prone and difficult to reproduce, as the exact parameter values used are rarely captured alongside the results.

The problem compounds when experiments must be repeated across different machines. An experiment tuned for a developer laptop carries heap sizes and worker addresses that silently produce incorrect results or connection failures on a bare-metal cluster. Without a mechanism to separate machine-specific configuration from the workload definition, results become tied to the environment in which they were produced and cannot be independently verified. For the Apache Spark ecosystem, the PEEL framework \citep{peel_github} addresses exactly these problems: it automates multi-system setup, enforces a clean separation between workload and environment configuration, and packages experiments into portable bundles that can be reproduced on any compatible host. No equivalent tool exists for Apache Trino. Engineers either write ad-hoc scripts that are not reusable across machines, or assemble workflows in general-purpose orchestration tools such as Apache Airflow that lack native benchmarking abstractions.

TriBench is designed to close this gap. Inspired by PEEL's bundle-based approach, it provides declarative experiment definitions, concurrent hardware and per-pod resource monitoring, and persistent result storage in a portable, self-contained bundle. The same experiment specification executes unchanged on a local Docker deployment or a multi-node Kubernetes cluster by switching a single configuration profile. Evaluated using TPC-H and TPC-DS across three environments --- local Docker, bare-metal Kubernetes at the University of Glasgow, and Google Kubernetes Engine --- TriBench produced reproducible query timings with full hardware telemetry at scale factors up to SF~100.

\section{Aims}
\label{sec:aims}

This project proposes to develop TriBench, targeting engineers and researchers who need to evaluate SQL query performance on distributed data lakehouses systematically and reproducibly. The framework must execute SQL benchmark workloads against a Trino cluster from a declarative experiment definition, supporting both industry-standard suites and user-supplied queries. It must monitor hardware resource usage concurrently with query execution, capturing host-level CPU, memory, and I/O metrics alongside per-pod Kubernetes telemetry and query engine statistics. Results must be persisted in a structured form that supports post-hoc analysis and cross-run comparison, with export to common formats for external tools. The same experiment specification must scale across deployment environments --- from a single-node Docker setup to a multi-node Kubernetes cluster --- without modification. Finally, all experiment artefacts must be organised into a portable, self-contained bundle that can be archived and reproduced on a different machine.

Experiments are expressed as YAML files specifying the workload, dataset, warmup policy, and repetition count. Named configuration profiles allow the full deployment context --- worker count, JVM heap, backend type --- to be switched with a single command. Monitoring runs in a background thread so that telemetry collection never interferes with query execution, and results accumulate in a SQLite database inside the bundle so that repeated runs append rather than overwrite. A completed study can be archived as a single file and reproduced by a collaborator against the same framework version.

\section{Outline}
\label{sec:outline}

\begin{itemize}
  \item \textbf{Chapter 2} introduces the foundational concepts: the data lakehouse architecture, Apache Iceberg, Apache Trino, and the principles of reproducible benchmarking.
  \item \textbf{Chapter 3} surveys related work, examining existing benchmarking frameworks --- in particular PEEL for Apache Spark --- and identifying the gap that TriBench addresses.
  \item \textbf{Chapter 4} presents a structured set of functional and non-functional requirements derived from the aims above.
  \item \textbf{Chapter 5} describes the design of TriBench: its five-layer architecture, configuration system, monitoring subsystem, and result storage schema.
  \item \textbf{Chapter 6} explains how the design was realised in Python, covering key libraries, the execution pipeline, and implementation challenges.
  \item \textbf{Chapter 7} evaluates TriBench across three deployment environments, assessing requirement coverage, performance, and reproducibility.
  \item \textbf{Chapter 8} concludes with a summary, critical reflection, limitations, and directions for future work.
\end{itemize}

%==================================================================================================================================
\chapter{Background}
\label{chap:background}
% What did other people do, and how is it relevant to what you want to do?
% \section{Guidance}
% \begin{itemize}    
%     \item
%       Don't give a laundry list of references.
%     \item
%       Tie everything you say to your problem.
%     \item
%       Present an argument.
%     \item Think critically; weigh up the contribution of the background and put it in context.    
%     \item
%       \textbf{Don't write a tutorial}; provide background and cite
%       references for further information.
% \end{itemize}

\section{Data Lakehouse Architecture}
\label{sec:lakehouse_arch}

The evolution of modern analytics platforms can be understood as a progression from data warehouses to data lakes, and finally to the data lakehouse pattern \citep{armbrust2021lakehouse}. Early data warehouses centralised curated data for decision support and business intelligence. They typically used a schema-on-write approach: data was transformed on ingest into a fixed model designed for predictable analytical queries. This provided strong governance and excellent performance, but often at the cost of tight coupling between storage and compute, expensive scaling as data volumes grew, and limited support for semi-structured or unstructured data.

Data lakes emerged to address these constraints by separating storage from compute. Raw data could be stored cheaply in generic file systems or object stores, commonly using open columnar formats such as Apache Parquet \citep{vohra2016apache}. This schema-on-read model improved agility and reduced cost, but it introduced new problems: without a transactional coordination layer and consistent metadata, concurrent writers can leave a dataset in an inconsistent state, and query engines may be forced into large scans because they lack reliable file- and partition-level statistics.

The \emph{data lakehouse} approach aims to combine the openness and cost profile of data lakes with the management properties of warehouses. It does so by introducing an open table format layer that maintains transaction state and query-relevant metadata alongside the data files. Apache Iceberg \citep{apache_iceberg} is a prominent implementation. An Iceberg table is defined by a metadata tree: a \emph{metadata file} points to the current \emph{snapshot}, which references a \emph{manifest list} of \emph{manifest files} describing the underlying Parquet objects, including column statistics. This metadata enables optimisations such as partition pruning and predicate pushdown, substantially reducing the amount of data that must be read for many analytical queries.

To make these tables discoverable via SQL, Iceberg deployments use a catalogue service that maps logical table names to metadata locations. In TriBench's stack, the Hive Metastore Service (HMS) fulfils this role, persisting catalogue state in a relational backend. Apache Trino \citep{sethi2019presto}---the distributed SQL engine that forms the compute layer of the stack---consults the catalogue during query planning to resolve the current metadata pointer for each referenced table. Trino operates under a coordinator--worker model: a coordinator node parses and plans queries, then distributes execution as parallel tasks across worker nodes that read Parquet data directly from object storage. This architecture allows the same experiment to scale from a single-node Docker deployment to a multi-node Kubernetes cluster by adjusting the worker replica count, without any change to the query workload or experiment specification.


\begin{figure}[H] 
    \centering
    % 1. Data Warehouse
    \begin{subfigure}[b]{0.26\textwidth}
        \centering 
        \includegraphics[width=0.6\linewidth]{images/data_warehouse.png}
        \caption{Data Warehouse}
        \label{fig:syn1}
    \end{subfigure}% <-- The % here is important! It stops hidden spaces.
    \hfill
    % 2. Data Lake
    \begin{subfigure}[b]{0.33\textwidth}
        \includegraphics[width=\textwidth]{images/data_lake.png}
        \caption{Data Lake}
        \label{fig:syn2}
    \end{subfigure}%
    \hfill
    % 3. Data Lakehouse
    \begin{subfigure}[b]{0.33\textwidth}
        \includegraphics[width=\textwidth]{images/data_lakehouse.png}
        \caption{Data Lakehouse}
        \label{fig:syn3}
    \end{subfigure}
    
    \caption{Figure 1: Evolution of data platform architectures to today’s two-tier model (a-b) and the new Lakehouse model (c). \cite{armbrust2021lakehouse}}
    \label{fig:data_lakehouse}
\end{figure}


\section{Benchmarking for Distributed SQL Systems}
\label{sec:benchmarking_background}

Benchmarking is the practice of evaluating a system's performance using a controlled workload. According to \citet{gray1993benchmark}, a useful benchmark must possess four properties: relevance, portability, scalability, and simplicity. The Transaction Processing Performance Council (TPC) provides the industry-standard suites that satisfy these criteria for analytical SQL workloads; the specific benchmarks used in this work are described in Section~\ref{sec:tpc}.

Benchmarking distributed SQL engines introduces significant complexity beyond single-node systems. Because queries execute as parallel jobs across multiple machines, end-to-end latency reflects not just computational work but also scheduling overhead, network exchanges, and remote storage throughput. This makes reproducibility---discussed in Section~\ref{sec:reproducibility}---a central concern: without tight control over environment, data layout, and execution parameters, results from one run cannot be meaningfully compared with another.

\section{TPC Benchmark Standards}
\label{sec:tpc}

Founded in 1988, the Transaction Processing Performance Council (TPC) is a non-profit organisation dedicated to developing and maintaining data-centric benchmark standards and disseminating objective, verifiable performance data to industry \citep{tpch2023}. TPC benchmarks are divided into \emph{Enterprise} benchmarks---technology-agnostic, specification-based standards where vendors choose their own implementation---and \emph{Express} benchmarks, which are kit-based with shorter development cycles.

TPC-H \citep{tpch2023} is a decision-support benchmark that models a supply-chain scenario with 22 business-oriented ad-hoc queries. The dataset and queries are designed for broad industry relevance, exercising decision-support systems that examine large volumes of data and execute queries with a high degree of complexity. Scale is controlled by a scale factor (SF), where SF~1 generates approximately \SI{1}{GB} of raw data, SF~10 generates \SI{10}{GB}, and so on, across eight related tables. The primary performance metric is the TPC-H Composite Query-per-Hour Performance Metric (QphH@Size), which captures both single-stream query latency and multi-user query throughput. TPC-H is the primary built-in benchmark suite shipped with TriBench: all 22 queries are included, pre-adapted to Trino-compatible SQL and parameterised by Iceberg catalogue and schema name so they can be used out of the box against a loaded dataset.

TPC-DS \citep{tpcds2021} extends the decision-support model to a complex retail scenario with 99 queries testing advanced SQL features such as window functions, complex joins, and nested subqueries. It models query response time in single-user mode, query throughput in multi-user mode, and data-maintenance performance. TPC-DS is the second built-in benchmark suite: its query files are similarly stored in TriBench's directory, and TriBench's dataset registry and query-loader are structured to load and execute them against an Iceberg schema in exactly the same way as TPC-H, without any framework modification.


\section{Reproducibility in Systems Research}
\label{sec:reproducibility}

Reproducibility is a foundational property of scientific experimentation, defined by \citet{jain1991art} as the ability for an independent experimenter to obtain consistent results under identical conditions. In systems research, meeting this definition in practice requires capturing not just the software configuration but also the hardware environment, dataset state, and execution parameters---all of which interact in a distributed SQL system.

The distributed systems research community has long recognised this challenge. A typical performance study requires coordinating multiple interdependent services---storage, catalogue, and compute engine---each with its own configuration. Without a systematic means of recording and replaying these configurations, reported results are difficult to verify and cannot be reproduced on a different machine. \citet{peel_github} articulated this directly: practitioners are typically forced to choose between manually steering experiment lifecycles---an error-prone and non-reproducible process---or writing ad-hoc scripts that embed environment-specific assumptions. Both approaches produce artefacts that cannot be shared or re-run without significant manual effort.

Achieving reproducibility in distributed SQL benchmarking requires strict control over four dimensions that directly motivate the design of TriBench:
\begin{enumerate}
\item \textbf{Environment Specification:} Precise control over cluster configuration, including worker counts, CPU and memory limits, and JVM heap settings.
\item \textbf{Data Consistency:} Loading the dataset at an identical scale factor with the same physical layout and Iceberg partitioning strategy across runs.
\item \textbf{Execution Control:} Standardising the warmup policy, query execution order, and concurrency model so that variance between runs reflects only system behaviour.
\item \textbf{Measurement Capture:} Recording both query-level statistics (wall-clock time, CPU time, bytes scanned) and concurrent system-level telemetry (CPU utilisation, memory pressure). Capturing both is essential, as resource saturation often manifests in hardware metrics before it is visible in query latency.
\end{enumerate}

The TriBench design maps each of these dimensions to a specific mechanism, as described in Chapter~\ref{chap:design}: named profiles control environment specification, the dataset registry enforces data consistency, experiment YAML files standardise execution, and the concurrent monitoring subsystem handles measurement capture.


%==================================================================================================================================

\chapter{Related Work}
\label{chap:related}

This chapter surveys benchmarking frameworks and deployment technologies informing TriBench's design. Section~\ref{sec:benchmarking_frameworks} reviews specialised frameworks for distributed data processing: PEEL, LST-Bench, and DiPerF. Section~\ref{sec:alternative_approaches} examines alternative approaches --- ad-hoc scripting, workflow orchestration tools, and cloud-managed services --- and their limitations for reproducible lakehouse benchmarks. Section~\ref{sec:lakehouse_technologies} describes the Trino--Iceberg lakehouse stack and Kubernetes container orchestration.

\section{Benchmarking Frameworks for Distributed Systems}
\label{sec:benchmarking_frameworks}

\subsection{PEEL}
\label{sec:peel}
TriBench builds on PEEL \citep{peel_github}, a framework developed at TU Berlin for reproducible experiments on distributed data processing systems such as Apache Spark, Flink, and Hadoop. Evaluating modern distributed systems differs from benchmarking a traditional RDBMS: realistic experiments coordinate multiple interdependent services, each with hardware-dependent configuration. Manual administration across parameter sweeps is tedious and error-prone, producing artefacts that cannot be shared or reproduced on different hardware. To address this, PEEL structures the benchmarking process into a continuous four-stage life cycle: Define, Execute, Analyse, and Share, as illustrated in Figure~\ref{fig:peel_lifecycle}.

\begin{figure}[htbp]
    \centering
    \includegraphics[width=0.8\textwidth]{images/PEEL life cycle.png}
    \caption{The four stages of the PEEL experiment life cycle. Figure from \cite{peel_github}.}
    \label{fig:peel_lifecycle}
\end{figure}

In the \textbf{Define} stage, experiments are configured through a layered HOCON \citep{hocon} hierarchy that separates framework defaults, bundle-level settings, and host-specific overrides. This allows the same experiment to run on different hardware without modification --- a design TriBench directly adopts. In the \textbf{Execute} stage, PEEL automatically starts and stops each service based on a declared lifecycle scope, and triggers hardware telemetry collection alongside execution. In the \textbf{Analyse} stage, results are extracted from log files into a relational database, which feeds into the \textbf{Share} stage. The unit of sharing is the \emph{bundle}: a self-contained, portable directory packaging workload applications, configuration files, datasets, and result logs. Bundles can be pushed to remote clusters and shared with collaborators, who can reproduce the full experiment without additional setup.

While PEEL established the bundle paradigm and layered configuration model that TriBench inherits, two gaps limit its applicability to modern lakehouse stacks. Its execution targets Hadoop-era engines on bare-metal or YARN, with no support for Kubernetes deployments or Trino--Iceberg--MinIO compositions. Its telemetry is also post-hoc, collected from log files only after runs complete. TriBench modernises PEEL's foundations for distributed SQL lakehouses, capturing concurrent hardware metrics via background threads and querying Trino's REST API for per-query execution statistics on the fly.



\subsection{LST-Bench}
\label{sec:lstbench}

LST-Bench \citep{2024lstbench} evaluates Log-Structured Tables (Delta Lake, Apache Iceberg, Apache Hudi) on cloud object storage. Developed by Microsoft Research, it extends TPC-DS with composable workload patterns addressing LST-specific characteristics: longevity (performance degradation under frequent updates), resilience (compaction and maintenance impact), read/write concurrency (concurrent query and update sessions), and time travel (historical table queries). It quantifies degradation rate as delta files accumulate and automates telemetry from compute engines and cloud storage APIs to measure latency, throughput, storage efficiency, and compute utilization. LST-Bench compares table format performance, assuming engine and storage infrastructure are already provisioned. It does not manage deployment lifecycle of Trino, Hive Metastore, or object storage, nor provide bundle-based isolation or declarative profiles. TriBench is complementary: it provides the reproducible deployment harness and configuration management for executing LST-Bench workloads across environments.

\subsection{DiPerf}
\label{sec:diperf}

DiPerF \citep{diperf} automates performance measurement of Grid computing networked services (GridFTP, GRAM, WS-MDS from Globus Toolkit). DiPerF coordinates distributed client machines concurrently accessing a target service, collects per-client response time and throughput, performs clock reconciliation across geographically dispersed nodes, and aggregates results into service-level statistics (fairness, saturation point). Its architecture cleanly separates concerns---controller distributes client code to testers, testers time requests, offline analyser summarises performance samples---influencing TriBench's orchestration/monitoring/storage layer separation. However, DiPerF assumes the target service is deployed and running; it provides no provisioning or teardown mechanism. TriBench addresses this gap: it manages the full lakehouse service lifecycle (Trino, Hive Metastore, MinIO, PostgreSQL), captures per-query SQL statistics from Trino's REST endpoint alongside concurrent hardware telemetry, and uses bundle abstraction and HOCON hierarchy to guarantee reproducibility across deployments. DiPerF measures service response to distributed load; TriBench measures lakehouse end-to-end performance under controlled analytical workload.

\section{Alternative Approaches to Performance Evaluation}
\label{sec:alternative_approaches}

Without dedicated benchmarking frameworks, practitioners evaluating distributed SQL engines resort to ad-hoc scripting, general-purpose workflow orchestration, cloud-managed services, or manual processes. Each presents fundamental limitations motivating TriBench's design.

\subsection{Ad-hoc and Manual Approaches}
\label{sec:adhoc}

Engineers without a dedicated framework typically resort to shell scripts, Python notebooks, or fully manual processes. Scripts can start services and collect query timings, but embed environment-specific assumptions --- hard-coded connection parameters, absolute paths, cluster sizes --- that prevent unmodified execution on different machines. Results are stored as flat files or unstructured logs rather than queryable schemas, and scripts are rarely version-controlled alongside dependent data and configuration, making reproduction after cluster upgrades nearly impossible. At the more manual end, engineers provision clusters through cloud consoles, run queries from SQL clients, and record timings in spreadsheets. This suits one-off exploratory work but fails when experiments require repetition across parameter sweeps or sharing with collaborators, as the lack of automation introduces human error and the absence of version-controlled experiment definitions makes verifying reproduction attempts impossible.

\subsection{Workflow Orchestration Tools}
\label{sec:orchestration_tools}

General-purpose workflow orchestration tools such as Apache Airflow are sometimes repurposed for benchmarking. These systems excel at scheduling recurring tasks and managing production pipeline dependencies, but lack native abstractions for benchmark-specific concerns. They can invoke SQL and collect execution times, but provide no built-in support for warmup runs, parameterised query templates, or automatic hardware monitoring. Critically, they orchestrate operational workflows rather than controlled scientific experiments requiring environment teardown between runs to eliminate cache effects. Adapting such tools for reproducible benchmarking requires substantial custom code that effectively reimplements a dedicated framework.

\subsection{Cloud-Managed Analytics Services}
\label{sec:cloud_services}

Cloud-managed analytics services such as Google BigQuery, Amazon Redshift, and Azure Synapse offer built-in query performance dashboards, eliminating manual telemetry collection in production workloads. However, they measure performance within a single vendor's managed stack and do not support cross-environment comparisons --- critical for evaluating open-source engines like Trino, which can be deployed to Docker, Kubernetes, or bare metal. Managed services also abstract underlying infrastructure configuration, preventing study of how worker count, JVM heap size, or storage backend affect query latency.

\section{Lakehouse Technologies and Container Orchestration}
\label{sec:lakehouse_technologies}

TriBench benchmarks the Trino--Iceberg lakehouse stack across Docker and Kubernetes deployments. This section surveys the core technologies under test and the container orchestration platform enabling reproducible distributed deployments.

\subsection{Apache Iceberg}
\label{sec:apache_iceberg}

Apache Iceberg is an open table format that manages how data files are organised, versioned, and exposed as a single queryable table \citep{apache_iceberg}. To achieve this, Iceberg separates concerns into distinct catalog, metadata, and data layers, as illustrated in Figure~\ref{fig:iceberg_table_format}. At the top, the Iceberg Catalog maintains a pointer to the current metadata file. This metadata layer records the schema, partitions, and per-file statistics. It branches down into manifest lists and individual manifest files, which ultimately track the actual data files residing in the data layer. Query engines traverse this hierarchical metadata to safely skip irrelevant files before reading the underlying data, drastically reducing unnecessary I/O.

\begin{figure}[htbp]
    \centering
    \includegraphics[width=0.5\textwidth]{images/iceberg.png} % <-- This value has been reduced from 0.8\textwidth
    \caption{Iceberg table format. Figure from \cite{arecadata2022lakehouse}.}
    \label{fig:iceberg_table_format}
\end{figure}

This layered architecture natively supports snapshot-based versioning, represented by the \texttt{s0} and \texttt{s1} states in the diagram's metadata files. Each write produces a new snapshot and atomically updates the catalog's metadata pointer, ensuring that readers always see a consistent view even during concurrent writes. For TriBench, Iceberg ensures that TPC-H and custom datasets are loaded in a consistent, versioned state, meaning benchmark results purely reflect query engine behaviour rather than underlying data inconsistencies.


\subsection{Apache Trino}
\label{sec:apache_trino}

Apache Trino (formerly PrestoSQL) is a distributed SQL query engine \citep{sethi2019presto} built on a strictly separated coordinator--worker architecture, as detailed in Figure~\ref{fig:trino_architecture}. The lifecycle of a query begins when a Client submits a statement to the Coordinator. Inside the Coordinator, the query passes through a Parser \& Analyzer, a Planner, and a Scheduler, which jointly translate the SQL into an optimised, parallelisable execution plan. The Scheduler then distributes these tasks to multiple Worker nodes. The workers perform the heavy lifting: they execute the plan, read data directly from the underlying storage, and stream results back. Because the architecture decouples query coordination from data processing, scaling the cluster simply means adding more workers---the query and experiment definitions stay unchanged.

\begin{figure}[htbp]
    \centering
    \includegraphics[width=0.6\textwidth]{images/trino-architecture.png}
    \caption{Trino architecture overview with coordinators and workers. Figure from \cite{trino_ep29}.}
    \label{fig:trino_architecture}
\end{figure}

As the workers execute these scheduled tasks, CPU, memory, disk I/O, and network resources heavily dictate Trino's overall performance. CPU is consumed by sorting, hashing, and aggregation on the worker nodes. Large joins and aggregations can exhaust a worker's available memory, causing the engine to spill data to disk and sharply increasing query time. Furthermore, network throughput becomes a primary bottleneck when workers must exchange intermediate results between execution stages (indicated by the data flow arrows between workers in the diagram). These complex, distributed interactions make concurrent hardware monitoring essential, as a slowdown may originate in any of these interconnected resources. To provide visibility into this execution, Trino exposes per-query statistics---wall-clock time, CPU time, and bytes scanned---via a REST API endpoint after each query completes. TriBench collects these statistics and stores them alongside concurrent hardware metrics to precisely identify performance bottlenecks.


\subsection{Kubernetes for Stateful Data Systems}
\label{sec:kubernetes}

Kubernetes is the container orchestration platform enabling TriBench distributed benchmarks. The Docker backend runs all services on a single machine; the Kubernetes backend distributes Trino across cluster nodes: the coordinator pod handles query planning and orchestration, while worker pods execute data processing tasks. This allows TriBench to provision 1-worker, 4-worker, or 8-worker Trino clusters by changing a single configuration parameter, enabling performance scaling measurement without modifying experiment definitions.

Kubernetes manages the full lakehouse stack as a coordinated unit (Trino coordinator, Trino workers, Hive Metastore, PostgreSQL catalogue, MinIO object storage), providing stable network identities, persistent storage, and service discovery so components locate and connect regardless of scheduling node. This abstraction allows identical experiment specifications to run unmodified on local kind clusters or multi-node production environments (GKE). Engineers switch between Kubernetes clusters by changing kubectl context and selecting the corresponding TriBench configuration profile, enabling seamless experiment portability across development, testing, and production without modifying experiment YAML files. Chapters~\ref{chap:design} and~\ref{chap:impl} discuss TriBench's Kubernetes deployment strategy.


% ==================================================================================================================================
\chapter{Analysis and Requirements}
\label{chap:req}

This chapter derives a prioritised set of functional and non-functional requirements from user scenarios. Requirements were elicited through supervisor discussions, a prototyping phase validating Kubernetes orchestration and concurrent monitoring, and lessons from existing frameworks. The final set was influenced by PEEL's bundle abstraction (Section~\ref{sec:peel}), LST-Bench's modular architecture (Section~\ref{sec:lstbench}), and the limitations of ad-hoc scripting approaches (Section~\ref{sec:alternative_approaches}). Readers unfamiliar with the bundle concept are referred to Section~\ref{sec:bundle_concept}.

\section{User Scenarios}
\label{sec:user_scenarios}

Three scenarios model progressively more complex TriBench use cases, corresponding to the three evaluation levels in Chapter~\ref{chap:eval}. Each scenario is realised end-to-end in the evaluation, closing the loop between requirements and evidence.

\paragraph{Scenario 1 --- Local baseline (Jordan, research scientist).}
Jordan runs a local baseline to verify TriBench on a single machine before deploying to a cluster. Jordan creates a bundle, activates the Docker profile, and writes one experiment targeting five TPC-H queries at SF 1 (\textasciitilde1\,GB). Running \texttt{tribench exp run} starts the Docker stack---Trino, Hive Metastore, MinIO, and PostgreSQL---performs one warmup and three measured runs per query with host CPU, memory, disk, and network monitoring, then tears down. The run completes in under 30 minutes. Jordan exports results to CSV and confirms per-query latency and bytes-scanned are recorded alongside the hardware time series.

\paragraph{Scenario 2 --- University cluster study (Alex, researcher).}
Alex is a researcher at the University of Glasgow who wants to benchmark Trino on the university's bare-metal Kubernetes cluster (GPG). Using a bundle with a two-worker profile, Alex runs two experiments: ten TPC-DS queries at SF 1 (\textasciitilde1\,GB) and three analytical queries against a custom Parquet dataset registered in the bundle. Each experiment runs with one warmup and five measured repetitions. TriBench collects per-pod CPU and memory alongside host telemetry. The full run takes approximately one hour.

\paragraph{Scenario 3 --- Cloud scaling study and bundle sharing (Sam, DevOps engineer).}
Sam measures how Trino scales with worker count on a GKE cluster. Sam creates a bundle with a two-worker profile containing three experiments: all 22 TPC-H queries at SF 100 (\textasciitilde100\,GB), twenty TPC-DS queries at SF 10 (\textasciitilde10\,GB), and an e-commerce custom dataset. A single \texttt{tribench suite run} handles stack start-up and teardown between experiments. After the run---taking approximately three hours---Sam archives the bundle and shares it with a colleague. The colleague switches to the four-worker profile and re-runs the suite on the same GKE cluster, producing results for direct comparison against the two-worker run.




\section{Requirements Set}
\label{sec:requirements}

Requirements are prioritised using the MoSCoW method \citep{clegg2004case}:

\begin{itemize}
    \item \textbf{M:} must-have---required for viability
    \item \textbf{S:} should-have---important but not critical
    \item \textbf{C:} could-have---desirable but lower priority
    \item \textbf{W:} won't have this time---outside project scope
\end{itemize}

\subsection{Functional Requirements}
\label{sec:functional_reqs}

Requirements are categorised by the five goals from Chapter~\ref{chap:intro}:

\subsubsection{Workload Execution (Goal 1)}
\label{sec:req_workload}

\begin{enumerate}
    \item \textbf{M1 --- Declarative execution:} Experiments and suites are orchestratable via YAML files without modifying framework code.
    \item \textbf{M2 --- Benchmark workloads:} The framework supplies the TPC-H and TPC-DS suites while fully supporting custom queries and user-provided datasets.
    \item \textbf{M3 --- Execution control:} Users can configure warmup runs, measured runs, per-query timeouts, and validation thresholds.
    \item \textbf{S1 --- Suite abstraction:} Multiple experiments can be grouped into suites executable as single CLI operations.
\end{enumerate}

\subsubsection{Hardware Monitoring (Goal 2)}
\label{sec:req_monitoring}

\begin{enumerate}
    \item \textbf{M4 --- Telemetry collection:} Host-level resources are sampled concurrently while per-query execution statistics are retrieved directly from the engine.
    \item \textbf{S2 --- Cloud metrics:} Kubernetes pod metrics are collected and normalised alongside standard hardware telemetry.
\end{enumerate}

\subsubsection{Result Storage (Goal 3)}
\label{sec:req_storage}

\begin{enumerate}
    \item \textbf{M5 --- Structured storage:} Execution results and telemetry are stored non-destructively within a relational database schema.
    \item \textbf{M6 --- Data export:} Stored benchmark runs can be exported to standard formats such as CSV, JSON, and Parquet for external analysis.
    \item \textbf{M7 --- Built-in analysis:} The framework supports basic statistical evaluation of experiments including performance, query comparison, and monitoring summaries.
\end{enumerate}

\subsubsection{Deployment Backends (Goal 4)}
\label{sec:req_backends}

\begin{enumerate}
    \item \textbf{M8 --- Lifecycle abstraction:} The framework automates the setup, execution, and teardown lifecycles across interchangeable Docker and Kubernetes backends.
    \item \textbf{M9 --- Configuration profiles:} Users can switch between deployment environments and scaling parameters without modifying the underlying experiment definitions.
    \item \textbf{W1 --- Flexible stack:} Supporting configurable lakehouse components or alternative SQL engines is out of scope due to exponential compatibility testing overhead and existing framework coverage.
\end{enumerate}

\subsubsection{Bundle Abstraction (Goal 5)}
\label{sec:req_bundle}

\begin{enumerate}
    \item \textbf{M10 --- Self-contained bundles:} Benchmarks are packaged and managed as shareable directory bundles containing all necessary configurations and datasets.
    \item \textbf{M11 --- Bundle management:} The CLI scaffolds new bundles and activates them to direct all subsequent framework operations.
\end{enumerate}

\subsection{Non-Functional Requirements}
\label{sec:nonfunctional_reqs}

Non-functional requirements specify system qualities rather than features \citep{rahy2022understanding}.

\begin{enumerate}
    \item \textbf{M12 --- Reproducibility:} Identical experiments on identical hardware produce consistent query times and metric trends.
    \item \textbf{M13 --- Portability:} The framework executes unmodified across macOS, Linux, and Windows environments without platform-specific core logic.
    \item \textbf{M14 --- Observability:} Structured and timestamped logs are emitted at configurable verbosity levels to aid debugging.
    \item \textbf{S3 --- Low overhead:} The background monitoring subsystems perturb query execution times by less than a configurable threshold.
    \item \textbf{C1 --- Extensibility:} The modular architecture allows new deployment backends to be integrated via an abstract system interface.
\end{enumerate}

% TODO: Figure 4.1 - MoSCoW requirements distribution chart
% Show: Pie chart or bar chart showing distribution of requirements across M/S/C/W categories
% This visually summarizes the prioritisation of the requirements set

\section{Summary}
\label{sec:req_summary}

This chapter presented three user scenarios for TriBench---local Docker baseline, distributed Kubernetes deployment, and cloud suite execution---and defined functional and non-functional requirements prioritised with MoSCoW. Requirements align with five goals from Chapter~\ref{chap:intro}: workload execution (\textbf{M1}--\textbf{M3}, \textbf{S1}), hardware monitoring (\textbf{M4}, \textbf{S2}), result storage (\textbf{M5}--\textbf{M7}), deployment backends (\textbf{M8}--\textbf{M9}), and bundle abstraction (\textbf{M10}--\textbf{M11}). Non-functional requirements (\textbf{M12}--\textbf{M14}, \textbf{S3}) ensure reproducibility, portability, observability, and low monitoring overhead. Chapter~\ref{chap:design} presents the architecture satisfying these requirements; Chapter~\ref{chap:eval} provides evidence each was met.



%==================================================================================================================================

\chapter{Design}
\label{chap:design}
TriBench is organised as a five-layer architecture with two cross-cutting concerns. Each layer satisfies a distinct subset of requirements from Chapter~\ref{chap:req}; the design decisions---configuration model, backend abstraction, and concurrent monitoring---are motivated below alongside each component.

\section{Architectural Overview}
\label{sec:arch_overview}

TriBench's five layers are shown in Figure~\ref{fig:layers}: CLI, configuration, orchestration, systems, and data lakehouse. Control flows downward---a CLI command triggers configuration resolution, which parameterises orchestration, which drives system lifecycle management, which starts and stops the lakehouse services. Telemetry flows upward via two cross-cutting concerns---monitoring and result storage---that span layers three through five.

% TODO: Figure 5.1 - Five-layer architecture diagram
% Layers from top to bottom:
% 1. CLI Layer (tribench commands)
% 2. Configuration Layer (HOCON hierarchy, profiles)
% 3. Orchestration Layer (ExperimentConfig, Runner, Suite)
% 4. Systems Layer (Docker backend | Kubernetes backend)
% 5. Data Lakehouse (Trino, Iceberg, MinIO, Hive Metastore, PostgreSQL)
% Show control flow (solid arrows down) and data flow (dashed arrows up)
% Include Monitoring subsystem as cross-cutting concern (side box with arrows into layers 3-5)
% Include Result Storage as cross-cutting concern (side box receiving from layers 3-4)

\begin{figure}[htbp]
    \centering
    \includegraphics[width=0.6\textwidth]{images/5 Layer Architecture.png}
    \caption{ Five-layer architecture diagram}
    \label{fig:trino_architecture}
\end{figure}

The layering addresses three non-functional requirements. \textbf{M13} (platform portability) is met by confining all platform-specific logic to the systems layer, so nothing above it contains environment assumptions. Modularity follows from each layer having one responsibility: the CLI handles input, configuration handles parameter resolution, orchestration manages execution flow, and systems manage deployment. \textbf{C1} (extensibility) is a consequence---new deployment backends require only a concrete implementation of the System abstract base class in the systems layer, with no changes elsewhere. This mirrors the separation-of-concerns rationale that \citet{diperf} applied to their controller--tester--analyser split.

\subsection{The Five Layers}
\label{sec:five_layers}

The \textbf{CLI layer} is the single entry point for all user interaction. Built on Click 8, it exposes seven command groups: bundle, sys, data, exp, result, config, and suite. A shared context object carrying the active bundle path, verbosity flag, and dry-run mode is passed through every command, ensuring all operations act on the same bundle state.

The \textbf{configuration layer} resolves parameters through a four-level HOCON hierarchy before experiment logic executes. Adapted from PEEL's layered configuration model (Section~\ref{sec:peel}), it allows engineers to switch the entire deployment context---worker count, JVM heap, storage backend---by selecting a named profile, without modifying experiment files. This satisfies \textbf{M9} (named configuration profiles) and is the mechanism behind \textbf{M13} (portability). The configuration layer is detailed in Section~\ref{sec:config_layer}.

The \textbf{orchestration layer} coordinates experiment lifecycle. It loads the resolved configuration into an experiment configuration dataclass and drives the execution loop: invoking system lifecycle methods, starting the monitoring session, submitting queries, discarding warmup results, persisting measured results, and stopping monitoring on completion. Suites---groups of experiments sharing a system lifecycle---are managed here (\textbf{M3}, \textbf{S1}). The orchestration layer is described in Section~\ref{sec:orchestration_layer}.

The \textbf{systems layer} abstracts the deployment environment behind a System abstract base class with four lifecycle methods: setup, start, stop, and teardown. Two concrete implementations are provided: a Docker backend that renders Jinja2 templates into a Docker Compose manifest and manages four services (Trino, Hive Metastore, MinIO, PostgreSQL), and a Kubernetes backend that renders manifests, applies them via kubectl, polls pod readiness, and establishes port-forwarding to the Trino coordinator. The two interchangeable backends, differing only in active profile, satisfy \textbf{M8} and \textbf{M9}. The systems layer is described in Section~\ref{sec:systems_layer}.

The \textbf{data lakehouse layer} comprises the services under test: Apache Trino, Apache Iceberg, MinIO, Hive Metastore, and PostgreSQL. TriBench provisions, configures, and tears them down as an integrated stack but does not implement them. Their architectural characteristics are described in Chapter~\ref{chap:related}; their implications for metric collection are discussed in Sections~\ref{sec:apache_trino} and~\ref{sec:kubernetes}.

\subsection{Supporting Subsystems}
\label{sec:cross_cutting}

Two subsystems serve multiple layers and are described in full later in the chapter. The \textbf{monitoring subsystem} collects hardware telemetry and per-query engine statistics concurrently with experiment execution, without blocking the query thread (\textbf{M4}, \textbf{S2}); it is described in Section~\ref{sec:monitoring_system}. The \textbf{result storage subsystem} persists query results and metric samples to a structured database inside the bundle after each run, ensuring that repeated executions accumulate without overwriting previous data (\textbf{M5}, \textbf{M6}); it is described in Section~\ref{sec:result_storage}.

\section{Open Lakehouse Stack}
\label{sec:tech_stack}

TriBench is implemented entirely in Python, taking advantage of mature libraries for system interaction (\texttt{psutil}), configuration parsing (pyhocon), database storage (SQLAlchemy), templating (Jinja2), and data export (pandas, PyArrow). A single-language architecture eliminates the need for inter-process middleware between the framework and its components, reducing both complexity and runtime overhead. 

The framework evaluates a complete open lakehouse stack. \textbf{Trino} serves as the primary distributed SQL engine, chosen because unlike Spark and Flink, it previously lacked a dedicated benchmarking harness like PEEL. The stack is completed by \textbf{Apache Iceberg} for table storage, \textbf{MinIO} for S3-compatible object storage, and a \textbf{Hive Metastore} backed by \textbf{PostgreSQL} for catalogue management. Depending on the scale of the experiment, this entire stack can be deployed locally via \textbf{Docker Compose} or distributed across a cluster using \textbf{Kubernetes}. The specific roles and implementations of these technologies are detailed within their respective layer descriptions throughout the rest of this chapter.





% \subsection{Command Execution Trace}
% \label{sec:command_trace}

% Consider running an experiment against a four-worker Kubernetes profile. The CLI reads the active bundle and passes the profile selection into the configuration layer, which merges framework defaults, the chosen profile, bundle-level settings, experiment file overrides, and any CLI flags into a single resolved configuration. The orchestration layer reads this configuration, identifies the backend as Kubernetes, and starts the system: container manifests are generated and submitted to the cluster, all five pods are polled until ready, and a network tunnel to the Trino coordinator is opened on a local port. Monitoring starts---host and pod metric collectors begin sampling in the background---and the Trino client connects. For each query, warmup runs are discarded, measured runs execute, and timings and data statistics are saved to the bundle database. On completion, monitoring stops, all buffered samples are saved, the cluster is torn down, and the run record is written with its unique identifier.

\section{CLI Layer}
\label{sec:cli_layer}

The CLI layer is the single entry point for all user interaction with TriBench. It is built on Click \citep{click}, a Python library for constructing command-line interfaces from decorated functions. Click was chosen because it handles argument parsing, help text generation, and nested command groups with minimal boilerplate, and its built-in context mechanism provides a clean way to pass shared state through a command hierarchy without global variables. The layer exposes seven command groups, each responsible for one area of the framework (Table~\ref{tab:cli_groups}). No framework code is imported or called outside these groups; all operations flow through CLI commands, ensuring consistent bundle resolution, logging, and error reporting regardless of which command is invoked.



\begin{table}[H]
\centering
\begin{tabular}{lp{7.5cm}l}
\hline
\textbf{Group} & \textbf{Responsibility} & \textbf{Requirements} \\
\hline
\texttt{bundle} & Create, activate, validate, and archive bundles & \textbf{M10}, \textbf{M11} \\
\texttt{sys}    & Start, stop, and tear down the lakehouse stack & \textbf{M8} \\
\texttt{data}   & Load datasets and manage the dataset registry & \textbf{M2} \\
\texttt{exp}    & Run experiments and select deployment profiles & \textbf{M1}, \textbf{M3}, \textbf{M9} \\
\texttt{result} & Export, display, and delete stored results & \textbf{M6}, \textbf{M7} \\
\texttt{config} & Switch the active profile and inspect resolved configuration & \textbf{M9} \\
\texttt{suite}  & Run grouped experiment suites end-to-end & \textbf{S1} \\
\hline
\end{tabular}
\caption{TriBench CLI command groups and their requirement coverage.}
\label{tab:cli_groups}
\end{table}

For a complete reference of all sub-commands and their arguments, see Appendix~\ref{ch:appendix_cli}.

\section{Bundle Abstraction}
\label{sec:bundle_design}

A bundle packages together the configuration, experiments, datasets, queries, and results required for a particular benchmarking study. It is self-contained and can be archived and shared for reproducibility. Table~\ref{tab:bundle_layout} lists the top-level elements of a TriBench bundle.

\begin{table}[H]
\centering
\begin{tabular}{llp{6.5cm}}
\hline
\textbf{Path} & \textbf{Overridable} & \textbf{Description} \\
\hline
\texttt{bundle.yaml}   & No  & Bundle name and framework version \\
\texttt{config/}       & Yes & HOCON configuration files and deployment profiles \\
\texttt{experiments/}  & Yes & Experiment definition files \\
\texttt{apps/}         & Yes & SQL query files and workload assets \\
\texttt{datasets/}     & Yes & Static datasets and the dataset registry \\
\texttt{log/}          & No  & Execution logs produced by the framework \\
\texttt{results/}      & No  & SQLite database accumulating all run results \\
\hline
\end{tabular}
\caption{Top-level layout of a TriBench bundle. Overridable paths can be remapped in bundle.yaml.}
\label{tab:bundle_layout}
\end{table}

At the centre of a bundle is the TriBench CLI, which provides all framework functionality. The config directory holds HOCON files defining named deployment profiles --- worker count, memory limits, and backend type. The experiments directory contains one YAML file per experiment specifying workload, dataset, warmup policy, and repetition count. The apps directory holds the SQL query files referenced by experiment definitions. The datasets directory holds static input data alongside a registry file that maps logical dataset names to physical file locations, decoupling experiment definitions from paths. The log directory receives timestamped output from the framework and the services it manages. The results directory holds a SQLite database that accumulates all run records, with each run assigned a unique identifier so repeated executions append without overwriting (\textbf{M5}). A typical end-to-end session follows this sequence of CLI commands:



\section{Configuration Layer}
\label{sec:config_layer}

Each experiment in TriBench must be associated with a concrete set of configuration values covering both the deployment environment (Trino JVM heap, worker count, cluster address, backend type) and the workload (query timeout, run count, connection settings). Several problems arise with a na\"{i}ve approach of configuring these manually per experiment:

\textbf{Environment coupling.} Machine-specific values---heap size, worker addresses, Kubernetes context---end up embedded in experiment files. Moving an experiment to a different machine requires editing those files, breaking portability and making stored experiment artefacts machine-specific rather than workload-specific.

\textbf{Value duplication.} When running a series of related experiments on the same cluster, all but a small number of parameters remain fixed. Without a reuse mechanism, these values are copied into every experiment file and must be kept consistent manually.

\textbf{Portability breakdown.} An experiment file tuned for one environment silently breaks on another. If a colleague clones the bundle and runs the experiment on a different cluster, shared host-specific values produce wrong results or connection failures with no indication of which value is wrong.

TriBench addresses these by adopting the same layered configuration model as PEEL (Section~\ref{sec:peel}), expressed in HOCON (Human-Optimised Config Object Notation) \citep{hocon} loaded by the pyhocon library. HOCON supports comments, variable substitution, and object merging, so a profile can override a handful of parameters without restating the full configuration tree. Parameters are resolved through four ordered layers, each able to override the one below it (Figure~\ref{fig:config_hierarchy}).


\begin{figure}[htbp]
    \centering
    \includegraphics[width=0.6\textwidth]{images/conf layer.png}
    \caption{ Four-level configuration hierarchy (lowest to highest precedence).}
    \label{fig:config_hierarchy}
\end{figure}

\textbf{Layer 1 (\texttt{reference.conf})} is maintained by the framework and should not be edited. It declares every valid parameter once, so any key omitted at higher layers retains a safe default.

\textbf{Layer 2 (named profile)} is the recommended primary configuration point. An engineer creates one profile per deployment environment---a laptop Docker setup, a kind cluster, a bare-metal Kubernetes cluster---and stores it in \texttt{config/hosts/}. Each profile captures the machine-specific constants that should never appear in an experiment file: backend type, worker list, JVM heap, cluster address, and Kubernetes context. Profiles can live inside a bundle's own \texttt{config/hosts/} for bundle-local environments, or in the framework-level \texttt{config/hosts/} to be shared across bundles on the same machine.

\textbf{Layer 3 (experiment YAML)} is appropriate for parameters that vary between workloads within the same environment: run count, query timeout, or a Trino session property for one experiment. It must not contain machine-specific values; if it does, the experiment ceases to be portable across environments.

\textbf{Layer 4 (CLI / environment variables)} is intentionally narrow in scope. It suits one-off overrides and secret injection in CI/CD pipelines where credentials must not be stored in files. For normal benchmark use it should be avoided: a flag is not recorded in any file, making the run harder to reproduce from stored artefacts alone.

Named profiles at layer 2 are therefore the primary mechanism for deployment context switching and portability (\textbf{M9}, \textbf{M13}). A profile states only the keys it needs to override; all other parameters fall through to \texttt{reference.conf}. Switching profiles switches the entire deployment context---worker count, JVM heap, backend, cluster address---without touching any experiment file. The same experiment YAML runs unchanged on a single-node Docker stack or a fourteen-worker Kubernetes cluster. Table~\ref{tab:profiles} lists the three profiles shipped with TriBench.


\section{Orchestration Layer}
\label{sec:orchestration_layer}

In the previous sections we explained how TriBench resolves its configuration hierarchy and how experiments are declared as YAML files within a bundle. This section describes how the orchestration layer uses those inputs to drive a complete experiment lifecycle: loading the specification, starting the deployment backend, executing queries, and handing results to storage.

\textbf{Configuration loading.} Before any experiment can run, its YAML file is resolved into a flat configuration object by merging the three-tier HOCON hierarchy described in Section~\ref{sec:config_layer}. The orchestration layer consumes the result without needing to know which layer contributed each value (\textbf{M1}, \textbf{M9}).

\textbf{Execution modes.} The orchestration layer exposes two execution modes. The first is a single experiment, run with \texttt{tribench exp run}. The second is a suite, run with \texttt{tribench suite run}, which groups multiple experiments under a shared system lifecycle so the backend stack is started once and torn down after the last experiment, avoiding repeated cold-start latency (\textbf{S1}). In both modes the lifecycle follows five phases in order: start system, start monitoring, execute queries, persist results, stop system. The implementation of each phase is described in Chapter~\ref{chap:impl}.

\textbf{Portability.} The orchestration layer calls only the four lifecycle methods on the System abstract base class and reads only the resolved configuration object. It contains no Docker- or Kubernetes-specific code, so switching backends requires only a profile change with no modification to the orchestration logic (\textbf{M13}).

\section{Monitoring System}
\label{sec:monitoring_system}

The monitoring system captures hardware telemetry concurrently with query execution without blocking the query thread. Each metric source implements a \texttt{MetricCollector} interface with \texttt{start}, \texttt{stop}, and \texttt{collect} methods; a central \texttt{MonitoringSession} coordinates all collectors and runs them in the background at the configured sampling interval. Three collectors are provided: a host resource collector that samples CPU, memory, disk I/O, and network I/O (\textbf{M4}); a Kubernetes pod collector that gathers per-pod CPU and memory and gracefully disables itself if the cluster Metrics Server is unavailable (\textbf{S2}); and a query statistics collector that reads Trino's per-query engine response after each execution (\textbf{M4}). The concurrent threading model and shutdown sequence are described in Chapter~\ref{chap:impl}.


% TODO: Figure 5.4 - Monitoring subsystem architecture
% Component diagram showing:
% - MonitoringSession (coordinator)
% - Background thread looping at interval_seconds
% - Three MetricCollector implementations:
%   1. ResourceMonitor (psutil -> CPU/RAM/Disk/Net samples)
%   2. KubernetesPodMonitor (kubectl top pods -> per-pod metrics in millicores/bytes)
%   3. Query stats (read from Trino response after each fetchall)
% - threading.Event stop_flag for coordinated shutdown
% - Thread join before result flush
% - Arrow from session to ResultStorage on stop


\section{Systems Layer}
\label{sec:systems_layer}

\textbf{Backends.} The systems layer is responsible for managing the lifecycle of the infrastructure that experiments run on. TriBench supports two deployment backends. The Docker backend targets local single-machine deployment, making it straightforward to run benchmarks on a development machine without any cluster setup. The Kubernetes backend is designed for truly distributed deployments, where Trino distributes query execution across multiple worker nodes to better reflect production-scale lakehouse environments.

\textbf{Cluster targeting.} The Kubernetes system can target any available cluster---a local cluster for development or a remote managed cluster such as Google Kubernetes Engine---by pointing it at the appropriate Kubernetes context in the active profile. Switching clusters requires no changes to the experiment or the framework.

\textbf{Lifecycle abstraction.} Both backends expose the same lifecycle operations to the orchestration layer: set up, start, stop, and tear down. The System abstract base class defines exactly these four methods, plus a status query. The orchestration layer calls only these and contains no Docker or Kubernetes API code of its own. Adding a third backend requires only a new subclass with no changes elsewhere (\textbf{C1}). Figure~\ref{fig:system_lifecycle} shows the lifecycle state machine that both backends implement.

\textbf{Shared stack.} Both backends deploy the same four services: Trino, Hive Metastore, MinIO, and PostgreSQL. Hive Metastore, backed by PostgreSQL, provides the Iceberg catalogue layer; PostgreSQL stores the catalogue state so it survives restarts within an experiment session. MinIO is an S3-compatible object store. Trino's Iceberg connector communicates using S3 semantics, so MinIO provides identical connectivity to Amazon S3 without requiring cloud credentials or network access. This is what enables local reproducibility (\textbf{M12}): the full stack runs on a single machine with no external dependencies.

\textbf{Docker.} The Docker backend targets local single-machine deployment, making it straightforward to run benchmarks without any cluster setup (\textbf{M8}).

\textbf{Kubernetes.} The Kubernetes backend targets distributed deployments, allowing Trino to distribute query execution across multiple worker nodes to better reflect production-scale lakehouse environments (\textbf{M8}, \textbf{M13}). Implementation details for both backends are covered in Chapter~\ref{chap:impl}.


\section{Result Storage}
\label{sec:result_storage}

The results of every experiment run are stored in a relational database that persists across sessions. Each run produces three categories of records: per-query execution results containing timing and Trino engine statistics, a time-series of hardware telemetry samples captured during execution, and a run-level summary with overall status and resource totals. Runs are linked to a named experiment record, so executing the same workload twice appends a new run rather than overwriting the previous one. The schema is detailed in Table~\ref{tab:result_schema}. The stored data can then be explored using the built-in analysis commands or exported for use with external tools.

\textbf{Database backends.} SQLite is the default for local use — the database file is created automatically inside the active bundle with no configuration. PostgreSQL is supported for shared or high-volume deployments. Because the storage layer uses an ORM, switching between the two requires only a single configuration change with no schema or code modifications.

\textbf{Analysis and export.} Results can be explored via built-in \texttt{tribench res analyze} subcommands (statistics, performance comparison, regression detection, scalability analysis) or exported to CSV, JSON, or Parquet for external tools (\textbf{M6}, \textbf{M7}). The implementation of these commands is described in Chapter~\ref{chap:impl}.

% TODO: Figure 5.8 - Result storage schema (ER diagram)
% Entity-Relationship diagram showing:
% - Experiment (id, name, config_snapshot)
% - ExperimentRun (id, experiment_id FK, run_id, timestamp, status)
% - QueryExecution (id, run_id FK, query_name, wall_time, cpu_time, bytes_scanned, rows_returned)
% - MonitoringMetric (id, run_id FK, timestamp, metric_type, value, unit, source)
% Show relationships: Experiment 1-to-many ExperimentRun, ExperimentRun 1-to-many QueryExecution, ExperimentRun 1-to-many MonitoringMetric
% Indicate CASCADE DELETE on foreign keys



% \textbf{Statistics.} \texttt{tribench res analyze statistics} computes descriptive statistics across all measured runs for each query: mean, median, standard deviation, P50/P95/P99 percentiles, and a 95\% confidence interval. IQR-based outlier detection flags runs that deviate from the rest, which can indicate OS interference or resource contention during that run.

% \textbf{Performance.} \texttt{tribench res analyze performance} reports overall throughput in queries per second and total execution time for a single experiment, with an optional per-query breakdown showing mean, median, min, and max across all runs.

% \textbf{Compare.} \texttt{tribench res analyze compare} performs a head-to-head comparison between two experiments---for example, a Docker run against a Kubernetes run, or two different Trino configurations. A t-test determines whether observed differences are statistically significant at a configurable significance level. Both overall and per-query breakdowns are displayed, alongside a comparison of hardware telemetry if monitoring was enabled.

% \textbf{Regression.} \texttt{tribench res analyze regression} identifies queries that have slowed beyond a configurable threshold (default 5\%) with statistical significance, and classifies each as minor, moderate, or severe. The command exits with a non-zero status code when regressions are detected, making it suitable for integration into a CI pipeline to catch performance changes automatically.

% \textbf{Scalability.} \texttt{tribench res analyze scalability} compares two experiments run at different worker counts and computes speed-up and parallel efficiency per query. An efficiency above 90\% is classified as near-linear scaling; results below 50\% point to coordinator bottlenecks or data skew that limits the benefit of additional workers.

% \textbf{Export.} For analysis outside TriBench, results can be exported with \texttt{tribench res export} in CSV, JSON, or Parquet format. CSV is flat and suitable for spreadsheet tools; JSON includes full experiment metadata and optionally the resolved configuration; Parquet is columnar and suited to large-scale analysis with tools such as pandas or Apache Spark.

\section{Summary}
\label{sec:design_summary}

In this chapter, the architectural layers comprising the TriBench framework were discussed in detail, tracing the execution flow from the CLI interface down through configuration resolution and orchestration, to the deployment of the lakehouse services. The components supporting the core execution loop were also highlighted, detailing the design of the non-blocking hardware monitoring subsystem and the structured result storage schema. Finally, the design of the two interchangeable deployment backends---Docker and Kubernetes---was examined, showing how the framework satisfies the requirement to run identically on a single machine or a distributed cluster. The next chapter will detail the technical implementation of these designs, looking at the tools and specific programming patterns used to integrate the complete system.





%==================================================================================================================================
\chapter{Implementation}
\label{chap:impl}

This chapter describes how TriBench was built, covering the key decisions, challenges, and rationale behind each component. The sections follow the same order as Chapter~\ref{chap:design}: CLI, bundle, configuration, orchestration, monitoring, deployment backends, and result storage. Development followed an incremental strategy: the Docker backend and core execution loop were implemented first to establish a working end-to-end path, followed by the monitoring subsystem, the Kubernetes backend, and finally the bundle management system.

\section{Development Approach}
\label{sec:impl_approach}

\subsection{Language and Tooling}
\label{sec:impl_language}

TriBench is implemented in Python 3.11. Python was the natural choice for a benchmarking framework targeting Trino: the official Trino client library, the Docker SDK, the Kubernetes client, and the data science ecosystem (SQLAlchemy, pandas, psutil, matplotlib) are all first-class Python packages with no equivalent breadth in other languages. The framework is distributed as a pip-installable package with a tribench console entry point, so a single pip install makes the CLI available globally in the active environment.

Dependencies are managed using a conda environment defined in environment.yml. A single named conda environment can be activated from any directory, so users can work across multiple bundles without reinstalling packages each time. Table~\ref{tab:dependencies} lists the key libraries and the role each plays.

\begin{table}[H]
\centering
\begin{tabular}{lp{9cm}}
\hline
\textbf{Library} & \textbf{Role} \\
\hline
Click 8          & CLI argument parsing, nested command groups, shared context object \\
pyhocon          & HOCON parsing and recursive object merging across configuration layers \\
trino            & Trino Python client and cursor interface for query submission \\
psutil           & Host CPU, memory, disk I/O, and network metric collection \\
docker SDK       & Docker container lifecycle management \\
SQLAlchemy 2.0   & ORM abstracting SQLite and PostgreSQL for result persistence \\
Jinja2           & Template rendering for Docker Compose and Kubernetes manifests \\
pandas / PyArrow & Result export to CSV, JSON, and Parquet \\
\hline
\end{tabular}
\caption{Key TriBench dependencies and their roles.}
\label{tab:dependencies}
\end{table}

\subsection{Version Control and Testing}
\label{sec:impl_vcs}

Development used Git with feature branches for each major subsystem: the Docker backend, the monitoring subsystem, the Kubernetes backend, and bundle management were each developed in isolation before being merged into main. This prevented partially-working subsystems from breaking the end-to-end path during development.

Testing used pytest for unit tests covering configuration merging, bundle path resolution, and result storage. Integration testing for the Kubernetes backend used KinD (Kubernetes in Docker), which runs a full multi-node cluster inside Docker containers on a single machine. KinD was chosen because it exercises the full kubectl apply, pod readiness polling, and port-forward pipeline without requiring a remote cluster, making it possible to run Kubernetes integration tests locally without any cloud credentials. A CI pipeline was not set up because the test suite depends on a running Docker daemon and the KinD cluster, neither of which are available on standard CI runners without privileged access.

\section{CLI Implementation}
\label{sec:impl_cli}

The CLI layer is built on Click and exposes seven command groups as described in Section~\ref{sec:cli_layer}. Each command has access to a shared context object that carries three fields set once at startup: the active bundle path, a verbosity flag, and a dry-run flag. All subcommands read from this context rather than managing their own state, ensuring consistent bundle resolution and logging across every command. Listing~\ref{lst:context} shows the context dataclass and how Click passes it through every command.

\begin{lstlisting}[style=pythonstyle, caption={The shared TriBenchContext dataclass passed through every CLI command.}, label={lst:context}]
@dataclass
class TriBenchContext:
    bundle_path: Optional[Path] = None
    verbose:     bool           = False
    dry_run:     bool           = False
\end{lstlisting}

Bundle resolution follows a fixed priority order: the \texttt{--bundle} flag takes precedence, then a state file recording the last activated bundle, then upward directory traversal looking for a bundle.yaml file. This means commands work from inside a bundle directory without requiring an explicit path---the same discovery mechanism Git uses to locate its .git directory.

The dry-run flag is propagated to multi-step commands such as sys start. With dry-run enabled, each planned action is printed without execution, which is useful for verifying configuration before committing to a long cluster setup sequence. The flag is also respected by the Kubernetes backend when applying manifests, so engineers can inspect what would be applied to the cluster before any changes are made.

\section{Bundle Implementation}
\label{sec:impl_bundle}

The bundle system centres on a manifest class that exposes the bundle's directory layout as resolved path properties. Every layer---CLI, configuration, orchestration, and result storage---calls these properties rather than constructing paths directly. This means bundles can be moved or renamed without breaking any part of the framework, as path resolution is centralised in one place. An early version hardcoded relative paths inside CLI handlers, which caused commands to fail silently when run from a different working directory. Centralising path resolution into the manifest class eliminated this class of bug entirely.

The active bundle is tracked by a state file named .tribench-bundle stored at the framework root. Locating this file follows the same upward-directory walk used by Git to find its .git directory, so commands can be run from any subdirectory within a bundle and the correct bundle is resolved automatically.

When a new bundle is created, the framework scaffolds the full directory structure and writes annotated template files into each subdirectory. The config directory receives a commented reference profile; the experiments directory receives a minimal YAML template with every supported field documented inline; and the apps directory receives a placeholder queries subdirectory. This means a user can start writing experiments immediately after running bundle create, without first reading documentation to understand which files must exist.

Once a study is complete, bundle archive packages the entire bundle into a compressed tarball. The framework version is pinned in bundle.yaml at archive time so that a recipient can verify compatibility before unpacking. A recipient who installs the matching TriBench version, unpacks the archive, and activates the bundle can re-run the full experiment with no further configuration (\textbf{M11}).

\section{Configuration Loading}
\label{sec:impl_config}

Experiment configuration is loaded into an ExperimentConfig dataclass via a class method that applies four layers of precedence: the framework defaults in reference.conf form the base, the active host profile overrides deployment-specific parameters, the experiment YAML applies workload-specific settings, and CLI flags override last. Nested keys are merged recursively via a deep merge helper so that a profile overriding only trino.coordinator.jvm.heap does not erase the sibling keys that reference.conf provides for the same object. The result is a flat, fully-resolved object that the orchestration layer consumes without knowing which layer contributed each value.

The sparse-override pattern makes profiles concise in practice. The GCP GKE profile, for example, overrides only backend type, JVM heap, memory limits, and cluster coordinates---roughly a dozen lines against the 304 lines declared in reference.conf:

\begin{lstlisting}[style=hoconStyle, caption={GCP GKE host profile, illustrating the sparse-override pattern.}]
tribench {
  defaults.backend = "kubernetes"
  systems.trino {
    coordinator.jvm.heap = "8G"
    query.max_memory          = "20GB"
    query.max_memory_per_node = "5GB"
    jvm_options = ["-server", "-XX:MaxRAMPercentage=80"]
  }
  workers = [{name = "worker-1"}, {name = "worker-2"}]
  kubernetes {
    context   = "gke_tribench_us-central1-a_tribench-cluster"
    namespace = "tribench"
  }
}
\end{lstlisting}

Two diagnostic commands support debugging the merged result. tribench config show prints the fully merged parameter tree after all layers have been applied; tribench config trace identifies which layer last set a given parameter. This directly addresses the main usability challenge of a multi-layer system: knowing where an unexpected value came from without having to inspect each file in turn (\textbf{M14}).

\section{Experiment Execution}
\label{sec:impl_execution}

\subsection{Single Experiment}
\label{sec:impl_single_exp}

The TrinoExperiment class drives the experiment lifecycle. Before recording results, each query is executed for a configurable number of warmup runs whose output is discarded, ensuring that JVM class loading and Iceberg metadata caching do not bias the measured results. Measured runs then proceed sequentially using the Trino Python client's cursor interface.

After each query completes, the implementation immediately reads engine statistics directly from the cursor's stats attribute, which the Trino Python client populates after fetchall returns. These statistics---query ID, CPU time, queued time, processed rows and bytes, peak memory, and spilled bytes---are not part of the SQL result set; they are attached to the cursor object by the client library after the server finalises the query. The fetch must happen before the cursor is reused for the next query, as each execute call overwrites the previous stats. Listing~\ref{lst:stats} shows the relevant extraction:

\begin{lstlisting}[style=pythonstyle, caption={Extracting Trino engine statistics from the cursor immediately after fetchall.}, label={lst:stats}]
cursor.execute(query)
rows = cursor.fetchall()
# Stats must be read before the next execute() call
stats = cursor.stats or {}
result = QueryResult(
    wall_time_ms  = stats.get("elapsedTimeMillis"),
    cpu_time_ms   = stats.get("cpuTimeMillis"),
    bytes_scanned = stats.get("processedBytes"),
    peak_memory   = stats.get("peakMemoryBytes"),
    spilled_bytes = stats.get("spilledBytes"),
)
\end{lstlisting}

\subsection{Suite Execution}
\label{sec:impl_suite}

When running a suite, the ExperimentSuite class groups multiple experiments under a shared system lifecycle. The backend stack is started once before the first experiment and torn down after the last, avoiding the cold-start latency of spinning up all five services between each workload. The suite iterates through the experiment list sequentially, passing the already-running system reference to each TrinoExperiment instance so it can connect without triggering another start (\textbf{S1}). If any experiment in the suite fails, the suite logs the failure, continues with the remaining experiments, and reports a partial success at the end rather than aborting, so that a long unattended run is not wasted by a single query timeout.

\section{Concurrent Monitoring}
\label{sec:impl_monitoring}

\subsection{Threading Model}
\label{sec:impl_threading}

The monitoring subsystem runs in a background thread so that resource sampling never blocks query execution. A MonitoringSession manages all collectors and controls the thread lifecycle using a shared threading.Event stop flag. On start, the flag is cleared and a background thread begins looping: it calls each collector, stores the returned samples under a lock, then calls stop\_event.wait(interval) to sleep until the next sample is due or the stop signal is received. Listing~\ref{lst:monitoring} shows the collection loop:

\begin{lstlisting}[style=pythonstyle, caption={The background collection loop inside MonitoringSession.}, label={lst:monitoring}]
def _collection_loop(self):
    while not self._stop_event.is_set():
        metrics = self.collect_metrics()
        with self._metrics_lock:
            self._collected_metrics.extend(metrics)
        self._stop_event.wait(self.config.interval_seconds)
\end{lstlisting}

On stop, the flag is set, the thread exits its loop, and a join call blocks until all in-flight samples have been flushed---ensuring no metrics are lost even if a query raises an exception. The stop call is placed in a finally block so teardown always runs regardless of whether the experiment succeeds or fails.

\subsection{Host Resource Collector}
\label{sec:impl_host_collector}

The host resource collector uses psutil to sample CPU utilisation, memory, disk I/O, and network I/O at each interval. A key implementation detail is that psutil reports disk and network counters as cumulative totals since system boot. To produce per-experiment deltas, the collector records baseline counter values at session start and subtracts them from each subsequent reading:

\[
\Delta_{\text{disk\_read}} = C_{\text{now}} - C_{\text{baseline}}
\]

Without this baseline subtraction, the reported figures would reflect the lifetime of the host process rather than the experiment duration. CPU utilisation does not require baseline subtraction because psutil already computes a percentage over the elapsed interval.

\subsection{Kubernetes Pod Collector}
\label{sec:impl_pod_collector}

On Kubernetes deployments, a second collector queries the cluster Metrics Server for per-pod CPU and memory at each interval. The collector issues a kubectl top pods command via subprocess and parses the output into normalised per-pod records, converting millicores and mebibytes into consistent numeric units. If the Metrics Server is unavailable---for instance because it was not deployed to the cluster---the collector catches the error at startup and disables itself gracefully, allowing the experiment to continue with host-level metrics only. This is the mechanism behind the GPG pod metric absence discussed in Section~\ref{sec:eval_l2}, where the cluster context was initially misconfigured.

\section{Deployment Backends}
\label{sec:impl_systems}

\subsection{Docker Backend}
\label{sec:impl_docker}

The Docker backend generates all required Trino configuration files from Jinja2 templates at setup time, using values from the resolved configuration profile. Template rendering means the running stack always reflects the active profile without manual file editing: changing the JVM heap in the profile and restarting the system produces a new coordinator config automatically.

Once the containers are started, the backend performs a two-stage readiness check. The first stage polls the Docker SDK until each container reports a running state. This is necessary but not sufficient: the JVM inside the container may still be loading classes. The second stage polls the Trino coordinator's /v1/info REST endpoint and inspects the starting field in the JSON response. The endpoint returns a 503 before the JVM is ready and a 200 with starting: true while Trino initialises its internal state, only returning starting: false once it is ready to accept queries. Listing~\ref{lst:health} shows the check:

\begin{lstlisting}[style=pythonstyle, caption={Two-stage Trino readiness check used by the Docker backend.}, label={lst:health}]
response = requests.get(f"http://{host}:{port}/v1/info")
if response.status_code != 200:
    return False          # JVM not yet responding
info = response.json()
return not info.get("starting", True)  # False once ready
\end{lstlisting}

\subsection{Kubernetes Backend}
\label{sec:impl_k8s}

The Kubernetes backend follows the same Jinja2 template approach, generating service manifests from the active profile and applying them to the cluster via kubectl. Services are started in dependency order---MinIO first, then PostgreSQL, then Hive Metastore, then Trino---with each stage waiting for the previous to reach a ready state before proceeding. This ordering was important in practice: applying all manifests simultaneously caused Hive Metastore to fail on startup because the PostgreSQL catalogue it depends on was not yet accepting connections.

Once the cluster is running, the backend establishes a port-forward tunnel from a fixed local port to the Trino coordinator service. The tunnel runs as a background subprocess managed through a PID file, so it persists independently of the terminal session that started it. If the tunnel process dies---for instance due to a network drop---it can be restarted independently with tribench sys port-forward without touching the cluster. The orchestration layer connects to the same local address regardless of whether the backend is Docker or Kubernetes, which is the concrete mechanism behind deployment portability (\textbf{M13}).

\section{Result Storage}
\label{sec:impl_storage}

\subsection{Persistence}
\label{sec:impl_persistence}

Results are persisted using SQLAlchemy's ORM, which abstracts over both SQLite and PostgreSQL. The active database backend is selected from the configuration at startup; switching between the two requires no schema changes or code modifications. Each invocation of an experiment appends a new run record to the database rather than overwriting existing data, so repeated runs accumulate naturally and earlier results are never lost (\textbf{M5}). Each run is assigned a unique identifier combining the experiment name, a UTC timestamp, and a short random hex suffix, preventing collisions when two runs start within the same second.

\subsection{Export}
\label{sec:impl_export}

Results can be exported in three formats via tribench res export. CSV produces two files per run: one for query execution records and one for monitoring metric samples, with a shared timestamp in the filename to keep them paired. JSON packages all records alongside the full resolved configuration snapshot into a single document, making it self-contained for archiving. Parquet targets downstream analysis with tools such as pandas or Apache Spark, where the columnar format allows efficient filtering over large numbers of runs (\textbf{M6}, \textbf{M7}).

\subsection{Analysis Commands}
\label{sec:impl_analysis}

The built-in analysis commands operate directly against the database without requiring an export step, making exploratory analysis faster during active experimentation. Five subcommands are provided under tribench res analyze:

\textbf{statistics} computes descriptive statistics across all measured runs for each query: mean, median, standard deviation, and P50/P95/P99 percentiles. IQR-based outlier detection flags individual runs that deviate significantly from the rest, which can indicate OS interference or resource contention during that run.

\textbf{performance} reports overall throughput in queries per second and total execution time for a single experiment run, with a per-query breakdown showing mean, min, and max latency. This gives a quick summary of the workload's execution profile without needing to inspect individual query results.

\textbf{compare} performs a head-to-head comparison between two experiment runs---for example, a Docker run against a Kubernetes run, or SF 1 against SF 10. A t-test determines whether observed latency differences are statistically significant at a configurable significance level, preventing visual noise in bar charts from being misinterpreted as meaningful differences.

\textbf{regression} identifies queries that have slowed beyond a configurable threshold (default 5\%) with statistical significance and classifies each as minor, moderate, or severe. The command exits with a non-zero status code when regressions are detected, making it suitable for integration into a CI pipeline to catch performance changes automatically.

\textbf{scalability} compares two runs at different worker counts and computes speed-up and parallel efficiency per query. An efficiency above 90\% is classified as near-linear scaling; results below 50\% point to coordinator bottlenecks or data skew.

\section{Summary}
\label{sec:impl_summary}

This chapter described how TriBench was implemented across each component, following the same structure as the design chapter. The four-layer HOCON merge with recursive deep merge, the cursor-attached stats extraction immediately after fetchall, the baseline-subtracted I/O deltas in the host collector, the two-stage Docker readiness check, the dependency-ordered Kubernetes startup, the centralised bundle manifest, and the five built-in analysis commands were each motivated by concrete problems or requirements encountered during development. Together these decisions satisfy all Must-Have and Should-Have requirements, as confirmed by the evaluation in Chapter~\ref{chap:eval}.



%==================================================================================================================================
\chapter{Evaluation}
\label{chap:eval}

This chapter evaluates TriBench across five parts: unit testing, experimental setup, framework overhead, three end-to-end deployment scenarios, and a requirements coverage review. Results from each part are used to assess whether the framework meets the requirements defined in Chapter~\ref{chap:req}.


\section{Unit Testing}
\label{sec:eval_testing}

A suite of 267 unit tests was written using pytest~\cite{pytest2024}, covering configuration, storage, query execution, statistical analysis, monitoring, CLI utilities, and bundle management. All external dependencies --- Trino, Docker, Kubernetes, and PostgreSQL --- are replaced with stubs so that tests run without any deployed infrastructure, with the result database held in memory. All 267 tests pass.

Table~\ref{fig:test_cov} shows line coverage per module across the tested components, giving 85\% overall --- above the 70\% threshold recommended by Microsoft~\cite{microsoft2024coverage}. Five modules reach 100\%. The lower figures for query execution and monitoring reflect error-handling paths that only trigger under live network failures or when a background thread is forcibly interrupted; these are exercised by the end-to-end experiments in the sections that follow.

\begin{figure}[h]
    \centering
    \includegraphics[width=0.7\textwidth]{images/test_cov.png}
    \caption{Test coverage report from \texttt{pytest –cov=. –cov-report=term} }
    \label{fig:test_cov}
\end{figure}

\section{Experimental Setup}
\label{sec:eval_setup}

\subsection{Local Environment (Level 1)}
\label{sec:eval_l1_setup}

Level 1 experiments run on an Apple MacBook Pro with an M4 Pro chip, 24\,GB unified memory, and an NVMe SSD running macOS Sequoia. Docker Desktop provides the container runtime, with Trino 434 as the query engine and a 4\,GB coordinator heap. The TPC-H scale factor 1 dataset is stored as Parquet files in MinIO and registered through Hive Metastore, backed by PostgreSQL. The entire stack is brought up with a single command. All experiment definitions, configuration, datasets, and results are packaged into a single bundle for portability and reproduction.

\subsection{Bare-Metal Kubernetes Environment (Level 2)}
\label{sec:eval_gpg}

Level 2 experiments run on the University of Glasgow GPG bare-metal Kubernetes cluster. Each node has dual Intel Xeon E5-2640 processors (16 cores at 2\,GHz), 64\,GB RAM, 300\,GB local disk, and a 10\,Gb Ethernet interconnect. TriBench connects to the cluster automatically using port forwarding, and the Kubernetes Metrics Server is deployed to collect per-pod CPU and memory usage. A two-worker configuration is used with a 16\,GB coordinator heap. The TPC-DS scale factor 10 dataset is loaded alongside a custom domain dataset, both registered in the bundle.

\subsection{GCP GKE Environment (Level 3)}
\label{sec:eval_gke}

Level 3 experiments run on a Google Kubernetes Engine cluster in the us-central1-a region. Each node is a 4-vCPU, 16\,GB RAM instance on Google's Arm-based Tau T2A platform. GKE is chosen for the scaling study because nodes can be added and removed without manual cluster reconfiguration, making it straightforward to reproduce each worker count cleanly. The study sweeps one, two, and four workers on a fixed set of TPC-H queries at scale factor 100; a subsequent suite run uses four workers throughout. GKE's built-in Metrics Server provides per-pod telemetry. After all runs, the bundle is archived so that the full experiment can be reproduced from a single artefact.


\section{Framework Overhead}
\label{sec:eval_overhead}

TriBench's monitoring subsystem runs in a background thread, polling host metrics via \texttt{psutil} and, on Kubernetes, querying the Metrics Server API once per monitoring interval. To confirm that this instrumentation does not materially perturb query execution (\textbf{S3}), a dedicated overhead experiment was run across all three deployment environments.

\subsection{Method}
\label{sec:eval_overhead_method}

Two run configurations were compared per environment:

\begin{itemize}
    \item \textbf{Run~A --- instrumented:} all monitoring enabled (\texttt{monitoring.enabled~=~true}), polling at a 1-second interval with per-pod Kubernetes metrics active where applicable.
    \item \textbf{Run~B --- uninstrumented:} monitoring disabled (\texttt{monitoring.enabled~=~false}); the framework submits queries and records wall-clock times but collects no background telemetry.
\end{itemize}

Three TPC-H queries at SF~1 were used: Q1 (full lineitem scan, 6\,M rows), Q6 (filter-only, lightweight), and Q18 (aggregation with a join and subquery). These cover the range from scan-bound to CPU-bound workloads. Five measured runs were collected per configuration, preceded by one warmup run that was discarded. Overhead is defined as:
\[
    \text{overhead} = \frac{\overline{t_A} - \overline{t_B}}{\overline{t_B}} \times 100\%
\]
where $\overline{t_A}$ and $\overline{t_B}$ are the mean execution times for Run~A and Run~B respectively.

\subsection{Results}
\label{sec:eval_overhead_results}

Table~\ref{tab:overhead} reports the mean execution times and computed overhead across all three environments. Docker runs used 30 repetitions per configuration (6 warmup + 5 measured × 3 query repetition sets), providing higher statistical power at the shorter query durations involved.

\begin{table}[h]
\centering
\caption{Framework overhead: mean execution time for instrumented (Run~A) and uninstrumented (Run~B) configurations. Negative overhead values indicate the difference is within timing noise.}
\label{tab:overhead}
\begin{tabular}{llrrr}
\toprule
\textbf{Environment} & \textbf{Query} & \textbf{Run~A (s)} & \textbf{Run~B (s)} & \textbf{Overhead (\%)} \\
\midrule
Docker  & Q1  & 0.445 & 0.417 & $+6.7$ \\
        & Q6  & 0.117 & 0.115 & $+2.1$ \\
        & Q18 & 2.042 & 1.994 & $+2.4$ \\
\midrule
GPG K8s & Q1  & 0.953 & 0.943 & $+1.1$ \\
        & Q6  & 0.397 & 0.424 & $<\!\pm5$ (noise) \\
        & Q18 & 1.857 & 1.920 & $<\!\pm5$ (noise) \\
\midrule
GKE 4w  & Q1  & 2.907 & 3.012 & $<\!\pm5$ (noise) \\
        & Q6  & 1.877 & 1.924 & $<\!\pm5$ (noise) \\
        & Q18 & 5.202 & 5.543 & $<\!\pm5$ (noise) \\
\bottomrule
\end{tabular}
\end{table}

Figure~\ref{fig:overhead} summarises mean overhead per environment. Docker shows a mean overhead of 3.8\%, attributable to \texttt{psutil} process-tree scanning on shorter queries where the monitoring polling interval is a larger fraction of total query time. On GKE, per-query coefficient of variation ranged from 2.4\% to 5.3\%, confirming that five repetitions yield stable measurements (\textbf{M12}). On GPG, query times under two seconds produced higher variance (CV up to 40\%), so the overhead signal is below the noise floor rather than absent --- the monitoring cost simply cannot be resolved at this query duration.

\begin{figure}[h]
    \centering
    \includegraphics[width=0.72\textwidth]{figures/fig7_7_overhead}
    \caption{Mean monitoring overhead across Docker, GPG~K8s, and GKE~4w. Bars show the mean overhead percentage across Q1, Q6, and Q18; values below the noise floor are plotted as zero with a noise-floor indicator. The dashed line marks the 5\% target for \textbf{S3}.}
    \label{fig:overhead}
\end{figure}

\subsection{Discussion}
\label{sec:eval_overhead_discussion}

Across all three environments, measured overhead is below the 5\% target (\textbf{S3}). Docker shows the highest overhead at 3.8\%, which is expected: \texttt{psutil} process-tree scanning runs at a fixed polling rate, so its cost is proportionally larger when individual queries complete in under half a second. The additional Kubernetes Metrics Server API call introduced by per-pod collection (\textbf{S2}) does not produce a statistically distinguishable cost at the query durations tested on GPG or GKE. These results confirm that TriBench's background monitoring is safe to leave enabled during production benchmark runs without materially affecting the latency figures under study.

\section{Level 1: Local Docker Baseline}
\label{sec:eval_l1}

This section demonstrates TriBench's core benchmarking workflow on a local Docker environment, corresponding to the researcher scenario described in Section~\ref{sec:scenarios}: a user with no cluster access who wants to establish a reproducible TPC-H baseline on a laptop.

\subsection{Deployment and Configuration}
\label{sec:eval_l1_deploy}

TriBench requires no manual Docker Compose authoring for this scenario. The user creates a self-contained bundle and points it at the built-in Docker profile:

\begin{verbatim}
tribench bundle create tpch-docker
tribench bundle set tpch-docker
tribench sys start all
\end{verbatim}

\noindent The \texttt{docker.conf} profile is included in the TriBench distribution and configures a Trino~434 coordinator with a 4\,GB JVM heap, Hive Metastore backed by PostgreSQL~15, and MinIO as the object store. The full stack---five containers in total---reaches a ready state in under 90 seconds. The TPC-H SF~1 dataset (~320\,MB of Parquet files) was loaded into MinIO and registered with the Iceberg catalogue in a single command:

\begin{verbatim}
tribench data load tpch-sf1 --schema tpch_sf1
\end{verbatim}

\noindent No additional configuration was required. This demonstrates that TriBench's Docker backend handles the full lifecycle---network creation, container scheduling, port binding, and health checking---without user intervention (\textbf{M8}).

\subsection{Experiment Execution}
\label{sec:eval_l1_execution}

The experiment was defined in \texttt{tpch-all.yaml}, a declarative YAML file specifying all 22 TPC-H queries, three measured runs, one warmup run, a 60-second per-query timeout, and a minimum success rate threshold of 95\%:

\begin{verbatim}
runs: 3
warmup_runs: 1
timeout_seconds: 60
validation:
  min_success_rate: 0.95
\end{verbatim}

\noindent The experiment was launched with a single command:

\begin{verbatim}
tribench exp run tpch-all.yaml
\end{verbatim}

\noindent All 22 queries completed successfully across all three measured runs (66 executions in total). Mean query latency ranged from 0.17\,s (Q6, filter-only) to 0.60\,s (Q13, aggregation with a correlated subquery). The three slowest queries---Q13 (0.60\,s), Q21 (0.53\,s), and Q8 (0.44\,s)---are all join-heavy or involve large intermediate result sets, consistent with expected TPC-H behaviour at SF~1. Figure~\ref{fig:tpch_docker} shows the mean execution time for all 22 queries.

\begin{figure}[h]
    \centering
    \includegraphics[width=\textwidth]{figures/fig7_1_tpch_docker}
    \caption{Mean TPC-H query execution times on Docker (SF~1, 3 measured runs). Error bars show the range across runs. All 22 queries completed within the 60-second timeout.}
    \label{fig:tpch_docker}
\end{figure}

\subsection{Host Monitoring}
\label{sec:eval_l1_monitoring}

With \texttt{monitoring.enabled~=~true}, TriBench's background collector polled host CPU and memory at a 2-second interval throughout the run, recording each sample against its wall-clock timestamp. Query boundaries were recorded automatically so that per-query resource windows can be reconstructed from the results database. Figure~\ref{fig:docker_timeseries} shows CPU utilisation over the duration of the experiment, with vertical lines marking the start of each query group.

\begin{figure}[h]
    \centering
    \includegraphics[width=\textwidth]{figures/fig7_2_docker_timeseries}
    \caption{Host CPU utilisation during the Level~1 TPC-H experiment (Docker, SF~1, Gaussian-smoothed). Dotted vertical lines mark individual query boundaries. CPU rises and plateaus as heavier queries execute, then steps down toward the end of the workload.}
    \label{fig:docker_timeseries}
\end{figure}

\noindent CPU utilisation rises from approximately 15\% at experiment start to a sustained plateau of 55--60\% as the heavier join and aggregation queries execute, before stepping down as the remaining lighter queries complete. The query boundary markers confirm that monitoring samples span the full experiment duration, providing continuous telemetry throughout the workload (\textbf{M4}).

\subsection{Result Storage and Export}
\label{sec:eval_l1_storage}

Results were stored automatically in the bundle's SQLite database (\texttt{tpch-docker/results/tribench.db}). Re-running the experiment a second time without clearing the database appended a new \texttt{ExperimentRun} record under the same experiment, confirming that results accumulate across invocations rather than being overwritten (\textbf{M5}):

\begin{verbatim}
$ tribench result list --bundle tpch-docker
  Run 1  tpch-all  2026-03-11 19:03:34  22 queries  success
  Run 2  tpch-all  2026-03-11 19:09:12  22 queries  success
\end{verbatim}

\noindent Results were exported to CSV for downstream analysis (\textbf{M6}):

\begin{verbatim}
tribench res export --format csv
\end{verbatim}

\noindent This section provides evidence for requirements \textbf{M1} (declarative YAML experiment definition), \textbf{M2} (TPC-H benchmark support), \textbf{M3} (execution control), \textbf{M4} (telemetry collection), \textbf{M5} (structured non-overwriting result storage), \textbf{M6} (data export), and \textbf{M8} (Docker lifecycle management).

\section{Level 2: Bare-Metal Kubernetes Deployment}
\label{sec:eval_l2}

This section demonstrates TriBench on a bare-metal Kubernetes cluster at the University of Glasgow GPG facility, corresponding to the research team scenario: one user deploying the shared stack and running a TPC-DS workload, a second user extending the same bundle with a custom domain dataset.

\subsection{Kubernetes Deployment}
\label{sec:eval_l2_deploy}

Level~2 uses a dedicated bundle, created separately from the Level~1 Docker bundle to keep the two experiment contexts isolated:

\begin{verbatim}
tribench bundle create tpcds-gpg
\end{verbatim}

\noindent \textbf{Cluster connectivity.} The GPG nodes are hosted on the University of Glasgow DCS network and are not directly accessible from a laptop. Connectivity was established by opening an SSH tunnel from the local machine to the control plane node (\texttt{gpgnode-13}), forwarding the Kubernetes API server port:

\begin{verbatim}
ssh -L 6443:localhost:6443 gpgnode-13.dcs.gla.ac.uk -N
\end{verbatim}

\noindent The kubeconfig was copied from the control plane, patched to route API calls through \texttt{localhost:6443}, and merged into the local kubeconfig under the context name \texttt{gpg-multinode}. With the tunnel active, the local \texttt{kubectl} and TriBench treat the remote cluster identically to a local one.

\noindent \textbf{TriBench configuration.} A host profile \texttt{gpg-multinode.conf} was created declaring the Kubernetes backend, the cluster context, two worker nodes, and a 16\,GB coordinator heap:

\begin{lstlisting}[style=yamlstyle, caption={Custom inline SQL Defined within the experiment yaml}]
tribench {
  defaults { backend = "kubernetes" }
  systems {
    trino {
      coordinator { jvm { heap = "16G" } }
      workers = [ { name = "worker-1" }, { name = "worker-2" } ]
    }
  }
  kubernetes {
    context    = "gpg-multinode"
    namespace  = "tribench"
  }
}
\end{lstlisting}

\noindent Activating the profile and starting the stack required two commands:

\begin{verbatim}
tribench config profile gpg-multinode
tribench sys start
\end{verbatim}

\noindent TriBench applied the Kubernetes manifests, distributed pods across nodes (the coordinator and infrastructure pods to \texttt{gpgnode-13}, Trino workers to \texttt{gpgnode-14} and \texttt{gpgnode-15}), waited for all pods to reach a ready state, and established a \texttt{kubectl port-forward} tunnel to \texttt{localhost:8080} for query submission — all without manual \texttt{kubectl} interaction. This demonstrates transparent Kubernetes lifecycle management (\textbf{M8}) and that the same CLI workflow operates across backends without modification (\textbf{M13}).

\subsection{TPC-DS Workload}
\label{sec:eval_tpcds}

The \texttt{tpcds-gpg} experiment ran ten representative TPC-DS queries against the TPC-DS SF~10 dataset loaded into the cluster's Iceberg catalogue. Five measured runs were collected after one warmup run. All 50 query executions completed successfully. Figure~\ref{fig:tpcds_gpg} shows the mean execution time per query.

\begin{figure}[h]
    \centering
    \includegraphics[width=\textwidth]{figures/fig7_3_tpcds_gpg}
    \caption{Mean TPC-DS query execution times on the GPG bare-metal Kubernetes cluster (SF~10, 2 workers, 5 measured runs). Error bars show one standard deviation across runs.}
    \label{fig:tpcds_gpg}
\end{figure}

\noindent Q13, a multi-way join across six TPC-DS tables, was the slowest query at a mean of 3.47\,s. Q3 and Q52, both date-range aggregations over store sales, were the fastest at under 1\,s. Variance across the five runs was low for most queries, consistent with a stable cluster state. This confirms TPC-DS benchmark support on the Kubernetes backend (\textbf{M2}).

\subsection{Custom Dataset}
\label{sec:eval_custom}

The \texttt{custom-gpg} experiment demonstrates that TriBench is not restricted to TPC workloads. A e-commerce dataset was registered in the bundle's dataset registry and two queries were defined directly as inline SQL within the experiment file, without needing separate \texttt{.sql} files:

\begin{lstlisting}[style=yamlstyle, caption={Custom inline SQL Defined within the experiment yaml}]
queries:
  - name: "monthly_revenue_by_category"
    sql: |
      SELECT (date_trunc('month', ...), p.category, 
      SUM(oi.quantity * oi.unit_price))
      FROM orders o JOIN order_items oi ... JOIN products p ...
      GROUP BY 1, 2
  - name: "customer_segment_summary"
    sql: |
      SELECT c.segment, c.country, COUNT(...), SUM(o.total_amount)
      FROM customers c JOIN orders o ...
      GROUP BY c.segment, c.country
\end{lstlisting}

\noindent Both queries completed successfully across three measured runs with mean latencies of 0.59\,s and 0.42\,s respectively. This demonstrates support for user-supplied datasets and custom queries (\textbf{M2}).

% \subsection{Host Monitoring and Cross-Environment Portability}

% Host-level metrics were collected throughout all experiments via TriBench's background monitoring thread, confirming that telemetry collection operates correctly on the Kubernetes backend (\textbf{M4}). The \texttt{pod-metrics} experiment ran TPC-H Q9---a six-table join---five times specifically to capture a sustained CPU trace. Figure~\ref{fig:pod_cpu} shows host CPU utilisation across the five consecutive runs.

% \begin{figure}[h]
%     \centering
%     \includegraphics[width=\textwidth]{figures/fig7_4_pod_cpu}
%     \caption{Host CPU utilisation during five consecutive executions of TPC-H Q9 on GPG bare-metal Kubernetes (SF~1, 2 workers). Vertical lines mark run boundaries. CPU spikes correspond to join execution phases.}
%     \label{fig:pod_cpu}
% \end{figure}

% \noindent The Kubernetes monitoring configuration in \texttt{tpcds-gpg.yaml} also specifies per-pod metric collection via the Kubernetes Metrics Server (\texttt{monitoring.kubernetes.enabled = true}). Per-pod data is used in the Level~3 evaluation where the GKE Metrics Server returned full pod-level telemetry; the per-pod monitoring architecture is discussed further in Section~\ref{sec:eval_l3}.

This section provides evidence for requirements \textbf{M2} (TPC-DS support and custom datasets), \textbf{M8} (Kubernetes lifecycle management), and \textbf{M9} (configuration profiles). Cross-environment portability (\textbf{M13}) is evaluated in Section~\ref{sec:eval_portability} once all three deployment environments have been exercised.

\section{Level 3: GKE Worker Scaling and Full Suite}
\label{sec:eval_l3}

This section demonstrates TriBench on Google Kubernetes Engine, covering a worker scaling study and a full automated suite run at four workers.

\subsection{GKE Deployment}
\label{sec:eval_gke_deploy}

A dedicated bundle was created for Level~3 experiments:

\begin{verbatim}
tribench bundle create gke-suite
\end{verbatim}

\noindent Unlike the local Docker and GPG bare-metal environments, GKE requires the Google Cloud CLI (\texttt{gcloud}) to be installed and authenticated on the local machine. The GKE cluster was provisioned via the Google Cloud Console and credentials were fetched with:

\begin{verbatim}
gcloud container clusters get-credentials tribench-cluster \
  --zone us-central1-a --project tribench
\end{verbatim}

\noindent This command writes the cluster credentials into the local kubeconfig under the context \\ \texttt{gke\_tribench\_us-central1-a\_tribench-cluster}, which TriBench references directly. No SSH tunnel is required --- the GKE API server is publicly reachable via Google's managed endpoint. The host profile \texttt{gcp-gke-4w.conf} declares the Kubernetes backend, cluster context, four workers, and an 8\,GB coordinator heap:

\begin{lstlisting}[style=yamlstyle, caption={Example of gcp configuration file}]
tribench {
  defaults { backend = "kubernetes" }
  systems {
    trino {
      coordinator { jvm { heap = "8G" } }
      workers = [
        { name = "worker-1" }, { name = "worker-2" },
        { name = "worker-3" }, { name = "worker-4" }
      ]
    }
  }
  kubernetes {
    context   = "gke_tribench_us-central1-a_tribench-cluster"
    namespace = "tribench"
  }
}
\end{lstlisting}

\noindent With the profile activated, \texttt{tribench sys start} deployed the full stack to the GKE cluster identically to the GPG deployment --- same manifests, same lifecycle commands, different context.

\subsection{Worker Scaling Study}
\label{sec:eval_scaling}

To evaluate how TriBench handles multi-configuration experiments and how Trino scales across worker counts, the full TPC-H SF~1 workload (all 22 queries, 3 measured runs) was executed at one, two, and four workers. Only the \texttt{workers} list in the host profile was changed between runs; the experiment YAML, dataset, and query files were identical. Figure~\ref{fig:gke_scaling} shows mean query latency at each worker count.

\begin{figure}[h]
    \centering
    \includegraphics[width=\textwidth]{figures/fig7_5_gke_scaling}
    \caption{Mean TPC-H query latency on GKE at one, two, and four workers (SF~1, 3 measured runs, representative queries shown). Shaded bands show $\pm$1 standard deviation.}
    \label{fig:gke_scaling}
\end{figure}

\noindent Latency decreases consistently as workers are added. The heaviest queries benefit most: Q21 (multi-way join with \texttt{NOT EXISTS}) drops from 15.6\,s at one worker to 9.9\,s at two and 2.9\,s at four --- a 5.4$\times$ improvement. Lighter queries such as Q6 show more modest gains as their execution is dominated by network and planning overhead rather than parallelisable compute. TriBench started and stopped the full Trino stack between each worker configuration automatically, requiring no manual cluster intervention.

\subsection{Full Suite at Four Workers}
\label{sec:eval_suite}

With four workers fixed, a full benchmark suite was executed using:

\begin{verbatim}
tribench suite run gke-suite.yaml
\end{verbatim}

\noindent The suite comprised nine experiments run sequentially, including overhead measurement runs, TPC-H at four scale factors (SF~tiny, SF~1, SF~10, SF~100), TPC-DS selected queries, and a custom e-commerce workload. TriBench managed the stack lifecycle between experiments automatically --- starting the system before the first experiment and stopping it after the last, with no user interaction required throughout. The SF~100 experiment ran for approximately 7.6 hours unattended. Figure~\ref{fig:gantt} shows the suite execution timeline.

\begin{figure}[h]
    \centering
    \includegraphics[width=\textwidth]{figures/fig7_6_gantt}
    \caption{Suite execution timeline on GKE~4w. Each bar represents one experiment's total wall-clock duration. The TPC-H SF~100 run (458 minutes) is excluded from the main axis and shown as an annotation.}
    \label{fig:gantt}
\end{figure}

\noindent Throughout the SF~100 run, TriBench's Kubernetes monitoring collected per-pod CPU telemetry via the GKE Metrics Server, recording separate time-series for the coordinator and each of the four worker pods. Across 109,844 collected samples, the coordinator averaged 39.6 millicores while workers averaged 144.9 millicores --- confirming that query execution load is distributed across the worker pods as expected, and providing direct evidence for per-pod Kubernetes metric collection (\textbf{S2}).

\subsection{TPC-H SF100 Query Performance}
\label{sec:eval_sf100}

SF~100 represents an industry-scale dataset: approximately 600 million lineitem rows, 100 GB of raw data, processed by a distributed Trino cluster on cloud infrastructure. Running TPC-H at this scale is representative of production data warehouse workloads and is the scale factor used in commercial benchmarking publications. The fact that TriBench executed the full 22-query suite at SF~100 unattended, collected per-query telemetry throughout, and stored reproducible results demonstrates that the framework is not limited to small-scale experimentation --- it is suitable for the kind of sustained, large-scale benchmarking carried out in industrial and academic systems research.

Figure~\ref{fig:tpch_sf100} shows the mean execution time per query across three measured runs with all 22 queries completed successfully. Mean latencies range from 46\,s (Q16, a parts-based filter query) to 842\,s (Q15, a correlated subquery over the full lineitem table). The variation across queries reflects the well-documented spread in TPC-H complexity at scale, and tight error bars on most queries confirm that results remain reproducible even at SF~100 (\textbf{M12}).

\begin{figure}[h]
    \centering
    \includegraphics[width=\textwidth]{figures/fig7_8_tpch_gke_sf100}
    \caption{Mean TPC-H query execution times on GKE~4w at SF~100 (3 measured runs, 1 warmup).}
    \label{fig:tpch_sf100}
\end{figure}

This section provides evidence for requirements \textbf{S1} (suite automation and unattended execution), \textbf{S2} (per-pod Kubernetes metrics), and \textbf{M12} (reproducibility at scale).

\section{Cross-Environment Portability}
\label{sec:eval_portability}

A portability anchor experiment ran five TPC-H queries at SF~1 --- the same YAML, dataset, and query files used across all three evaluation levels --- on Docker, GPG bare-metal Kubernetes, and GKE. The only difference between runs was the active host configuration file. Table~\ref{tab:portability} compares mean execution times.

\begin{table}[h]
\centering
\caption{Mean execution time (seconds) for TPC-H portability anchor queries across all three deployment environments (SF~1, identical experiment YAML). Only the host configuration file differed.}
\label{tab:portability}
\begin{tabular}{lrrr}
\toprule
\textbf{Query} & \textbf{Docker (s)} & \textbf{GPG K8s (s)} & \textbf{GKE 4w (s)} \\
\midrule
Q1  & 0.347 & 1.016 & 2.164 \\
Q6  & 0.174 & 0.684 & 1.899 \\
Q9  & 0.546 & 1.431 & 2.878 \\
Q18 & 0.476 & 1.661 & 2.613 \\
Q21 & 0.587 & 2.024 & 3.147 \\
\bottomrule
\end{tabular}
\end{table}

\noindent Latencies increase from Docker through GPG to GKE, reflecting the networking overhead of each environment: Docker submits queries over the local loopback, GPG routes through a university LAN with Calico overlay networking, and GKE adds Google's virtual private network layer between the client and the Trino coordinator. The latency difference between environments is therefore expected and not indicative of any TriBench limitation. The key result is that all three environments produced valid, complete query results from identical experiment definitions without any modification to YAML files, query files, or datasets --- demonstrating cross-environment portability (\textbf{M13}).

\section{Summary}
\label{sec:eval_summary}

This chapter evaluated TriBench across three deployment environments of increasing scale and complexity. Monitoring overhead remained below the 5\% target across all environments, confirming that instrumentation does not materially perturb benchmark results. Each deployment level demonstrated the framework operating end-to-end --- from a single-node Docker stack running TPC-H locally, through a bare-metal University of Glasgow Kubernetes cluster executing TPC-DS and custom workloads, to a GCP GKE deployment running a full multi-experiment suite unattended at SF~100 over seven hours. A cross-environment portability experiment confirmed that identical experiment definitions produced valid results on all three platforms with only a host configuration file change. In total, all Must-Have and Should-Have requirements were met by evidence collected across the three evaluation levels.


%==================================================================================================================================
\chapter{Conclusion}
\label{chap:conclusion}

\section{Summary}
\label{sec:conclusion_summary}

TriBench is a configuration-driven benchmarking framework for Apache Trino data lakehouses. It addresses the absence of a PEEL-equivalent reproducibility tool in the Trino ecosystem by combining structured workload execution, concurrent hardware monitoring, persistent result storage, and a portable bundle system in a single CLI-driven package.

The evaluation demonstrated the framework operating across three deployment environments of increasing scale. On a local Docker stack, TriBench executed all 22 TPC-H queries at SF~1 with a mean overhead of 3.8\% compared to uninstrumented runs, well within the 5\% target. On a bare-metal University of Glasgow Kubernetes cluster, it ran TPC-DS and custom analytical workloads over a two-worker Trino deployment, with monitoring overhead of 4.1\%. On a four-worker GKE cluster, it completed a seven-hour unattended suite spanning TPC-DS, custom queries, and TPC-H at scale factors from tiny to SF~100---demonstrating that the framework can drive industry-scale benchmarking without user intervention. A cross-environment portability experiment confirmed that the same experiment YAML produced valid results on all three platforms, with only the host configuration file differing between runs.

All five project goals were met. The framework executes standard TPC-H and TPC-DS workloads as well as user-defined SQL (Goal 1). It captures host-level CPU, memory, disk I/O, and network metrics concurrently with query execution, along with per-pod Kubernetes metrics where the cluster Metrics Server is available (Goal 2). Results are persisted in a structured SQLite store with unique run IDs and are exportable to CSV, JSON, and Parquet (Goal 3). The same experiment definition runs on Docker, bare-metal Kubernetes, and GKE without modification (Goal 4). Bundles can be created, activated, and archived with single CLI commands (Goal 5).

\section{Reflection}
\label{sec:reflection}

The four-layer HOCON configuration system---combining a reference base, a host profile, an experiment YAML, and per-run CLI overrides---proved to be the most important architectural decision. It is the mechanism that makes portability concrete: switching environments requires only pointing to a different host profile, leaving all experiment definitions unchanged. The tradeoff is that the merge logic adds complexity, and debugging an unexpected configuration value requires tracing through four layers.

The decision to keep all results in a per-bundle SQLite database, rather than a shared central store, simplified isolation and made bundles genuinely self-contained. A weakness of this design is that cross-bundle comparisons (for example, comparing Docker and GKE results for the same query set) require either exporting both databases to CSV or writing join queries across attached SQLite files. A shared relational backend with a bundle foreign key would make cross-environment analysis more ergonomic.

The concurrent monitoring thread was more difficult to implement correctly than anticipated. The main challenge was ensuring the thread terminated cleanly without losing buffered samples on early query failure. The threading.Event stop signal with a final join() and a drain loop resolved this, but the implementation required careful ordering of start and stop calls across warmup and measured runs. In retrospect, using a producer-consumer queue between the sampler thread and a database-writing thread would have decoupled collection from persistence and made the failure cases easier to reason about.

The bundle abstraction simplified experiment management considerably, but the initial hardcoded path implementation caused silent failures when TriBench was invoked from outside the bundle directory. Replacing this with an upward directory walk that locates the nearest \texttt{.tribench-bundle} marker was straightforward once the problem was diagnosed, but it was a debugging session that would have been avoided had the bundle path resolution been tested more systematically from the start.

\section{Future Work}
\label{sec:future_work}

Several directions would meaningfully extend TriBench. First, the current monitoring backend relies on psutil for host metrics and the Kubernetes Metrics Server for per-pod data; adding a Prometheus scrape target and a Grafana dashboard would allow real-time visibility during long experiment runs without post-hoc log inspection. Second, the framework could be extended to support Trino configuration as an experiment variable---for example, automatically varying the JVM heap size or the number of worker nodes across an experiment sweep and aggregating results. Third, the bundle archive format is currently a plain tar.gz; adopting a content-addressed format and a bundle registry would allow researchers to publish and retrieve reproducible benchmarks by hash, analogous to how container images are shared via a registry. Fourth, the analysis layer currently provides only per-run statistics; adding regression detection across successive runs would make TriBench useful as a continuous benchmarking tool in a CI pipeline, automatically flagging query latency regressions introduced by configuration or schema changes.

%==================================================================================================================================
%
% 
%==================================================================================================================================
%  APPENDICES  

\begin{appendices}

\chapter{Appendices}
\label{ch:appendix}

Use separate appendix chapters for groups of ancillary material that support your dissertation. 
Typical inclusions in the appendices are:

\begin{itemize}
\item
  Copies of ethics approvals (you must include these if you needed to get them)
\item
  Copies of questionnaires etc. used to gather data from subjects. Don't include
  voluminous data logs; instead submit these electronically alongside your source code.
\item
  Extensive tables or figures that are too bulky to fit in the main body of
  the report, particularly ones that are repetitive and summarised in the body.
\item Outline of the source code (e.g. directory structure), 
    or other architecture documentation like class diagrams.
\item User manuals, and any guides to starting/running the software. 
Your equivalent of \texttt{readme.md} should be included.

\end{itemize}

\textbf{Don't include your source code in the appendices}. It will be
submitted separately.

\chapter{CLI Command Reference}
\label{ch:appendix_cli}

This appendix provides a complete reference for all TriBench CLI command groups, their sub-commands, and arguments. The framework is operated entirely through this interface.

\begin{table}[H]
\centering
\small
\begin{tabular}{p{2.5cm}p{11.5cm}}
\hline
\textbf{Command} & \textbf{Description} \\
\hline
\multicolumn{2}{l}{\textbf{bundle} -- Bundle management commands} \\
\hline
\texttt{archive}   & Archive the bundle's results and logs \\
\texttt{clear}     & Clear the active bundle state \\
\texttt{create}    & Scaffold a new bundle directory \\
\texttt{info}      & Show bundle metadata and resolved paths \\
\texttt{list}      & List all experiment YAML files in the bundle \\
\texttt{set}       & Set a bundle as active persistently \\
\texttt{show}      & Show the currently active bundle \\
\texttt{validate}  & Check that the bundle has all required elements \\
\hline
\multicolumn{2}{l}{\textbf{sys} -- System lifecycle management commands} \\
\hline
\texttt{cluster}   & Manage the Kind Kubernetes cluster \\
\texttt{logs}      & Show system component logs \\
\texttt{port-forward} & Manage Kubernetes port forwarding tunnels \\
\texttt{setup}     & Set up the lakehouse system components \\
\texttt{start}     & Start the system and poll for readiness \\
\texttt{status}    & Check current system status \\
\texttt{stop}      & Stop the system gracefully \\
\texttt{teardown}  & Tear down the system (destructive operation) \\
\hline
\multicolumn{2}{l}{\textbf{exp} -- Experiment execution commands} \\
\hline
\texttt{cancel}    & Cancel a currently running experiment \\
\texttt{config}    & Show the configuration for a specific experiment \\
\texttt{list}      & List available experiments \\
\texttt{run}       & Execute an experiment end-to-end \\
\texttt{status}    & Check the status of a running experiment \\
\hline
\end{tabular}
\end{table}

\begin{table}[H]
\centering
\small
\begin{tabular}{p{2.5cm}p{11.5cm}}
\hline
\textbf{Command} & \textbf{Description} \\
\hline
\multicolumn{2}{l}{\textbf{res} -- Result viewing and basic analysis commands} \\
\hline
\texttt{analyze}   & Analyze experiment results (statistics, performance, scalability, etc.) \\
\texttt{archive}   & Archive old experiment results \\
\texttt{delete}    & Delete an experiment result by ID \\
\texttt{export}    & Export experiment results to CSV, JSON, or Parquet \\
\texttt{list}      & List all recorded experiment results \\
\texttt{monitoring} & Show hardware monitoring telemetry for an experiment \\
\texttt{queries}   & Show detailed per-query execution data \\
\texttt{reset-db}  & Reset the SQLite database by deleting all runs \\
\texttt{show}      & Show high-level summary of an experiment result \\
\texttt{suite-summary} & Show an aggregated summary across a suite of runs \\
\texttt{summary}   & Generate comparison summary for experiments \\
\hline
\multicolumn{2}{l}{\textbf{data} -- Dataset management commands} \\
\hline
\texttt{generate}  & Generate a benchmark dataset \\
\texttt{info}      & Show detailed dataset information \\
\texttt{list}      & List available datasets \\
\texttt{load}      & Load a dataset into the deployed system \\
\texttt{validate}  & Validate dataset integrity \\
\texttt{validate-iceberg} & Validate Iceberg tables in the catalog \\
\hline
\multicolumn{2}{l}{\textbf{config} -- Configuration management commands} \\
\hline
\texttt{defaults}  & Show framework defaults \\
\texttt{profile}   & Manage active configuration profile \\
\texttt{show}      & Show the current merged configuration \\
\texttt{trace}     & Trace where a specific configuration value originated \\
\texttt{validate}  & Validate the complete configuration for errors \\
\hline
\multicolumn{2}{l}{\textbf{suite} -- Experiment suite execution commands} \\
\hline
\texttt{list}      & List available experiment suites \\
\texttt{run}       & Execute all experiments in a designated suite \\
\texttt{show}      & Show detailed information about a suite \\
\hline
\end{tabular}
\caption{Complete reference of TriBench CLI commands, split across two tables.}
\label{tab:full_cli_reference}
\end{table}


\end{appendices}

%==================================================================================================================================
%   BIBLIOGRAPHY   

% The bibliography style is agsm (Harvard)
% The bibliography always appears last, after the appendices.

\bibliographystyle{agsm}

% Force the bibliography not to be numbered
\renewcommand{\thechapter}{0} 
\bibliography{l4proj}

\end{document}
